\subsection{CENTRE OF THE ENVELOPING ALGEBRA}

In this number, we choose a basis B of R. Let \(R_{+}\) be the set of positive roots relative to B. Let \(\rho = \frac{1}{2} \sum_{\alpha \in R_{+}} \alpha\) , and \(\delta\) the automorphism of the algebra \(\mathbf{S}(\mathfrak{h})\) that takes every \(x \in h\) to \(x - \rho(x)\) , and hence the polynomial function p on \(h^{*}\) to the function \(\lambda \mapsto p(\lambda - \rho)\) .

THEOREM 2. Let U be the enveloping algebra of g, Z its centre, V ⊂ U the enveloping algebra of h (identified with S(h)), U⁰ the commutant of V in U, \(\varphi\) the Harish-Chandra homomorphism (§6, no. 4) from \(U^{0}\) to V relative to B. Let \(\mathbf{S}(\mathfrak{h})^{\mathrm{W}}\) be the set of elements of \(\mathbf{S}(\mathfrak{h})\) invariant under the action of W. Then \((\delta\circ\varphi)|Z\) is an isomorphism from Z to \(\mathbf{S}(\mathfrak{h})^{\mathrm{W}}\) , independent of the choice of B.

\begin{enumerate}
\def\labelenumi{\alph{enumi})}
\tightlist
\item
  Let \(P_{++}\) be the set of dominant weights of R, \(w \in W\) , \(\lambda \in P_{++}\) , \(\mu = w\lambda\) . Then \(Z(\mu - \rho)\) is isomorphic to a submodule of \(Z(\lambda - \rho)\) (§6, no. 3, Cor. 2 of Prop. 6), and \(\varphi(u)(\lambda - \rho) = \varphi(u)(\mu - \rho)\) for all \(u \in Z\) (§6, no. 4, Prop. 7). Thus, the polynomial functions \((\delta \circ \varphi)(u)\) and \((\delta \circ \varphi)(u) \circ w\) on \(h^{*}\) coincide on \(P_{++}\) . But \(P_{++}\) is dense in \(h^{*}\) in the Zariski topology: this can be seen by identifying \(h^{*}\) with \(k^{B}\) by means of the basis consisting of the fundamental weights \(\varpi_{\alpha}\) , and by applying Prop. 9 of Algebra, Chap. IV, §2, no. 3. Hence
\end{enumerate}

\[
(\delta \circ \varphi) (u) = (\delta \circ \varphi) (u) \circ w,
\]

which proves that \((\delta \circ \varphi)(\mathbf{Z}) \subset \mathbf{S}(\mathfrak{h})^{\mathrm{W}}\) .

\begin{enumerate}
\def\labelenumi{\alph{enumi})}
\setcounter{enumi}{1}
\tightlist
\item
  Let \(\eta\) be the isomorphism from \(\mathrm{I}(\mathfrak{g})\) to \(\mathbf{S}(\mathfrak{h})^{\mathrm{W}}\) defined in no. 3, Cor. 2 of Th. 1. Consider the canonical isomorphism from the g-module U to the g-module \(\mathbf{S}(\mathfrak{g})\) (Chap. I, §2, no. 8), and let \(\theta\) be its restriction to Z. Then \(\theta(\mathrm{Z}) = \mathrm{I}(\mathfrak{g})\) . Let z be an element of Z with filtration \(\leq f\) in U.
\end{enumerate}

\[
\begin{array}{c c c} \mathrm{Z} & \stackrel {{\theta}} {{\longrightarrow}} & \mathrm{I} (\mathfrak {g}) \\ \varphi \Big \downarrow & & \Big \downarrow^ {\eta} \\ \mathbf {S} (\mathfrak {h}) & \stackrel {{\delta}} {{\longrightarrow}} & \mathbf {S} (\mathfrak {h}). \end{array}
\]

Introduce the notations of §6, no. 4, and put

\[
z = \sum_ {\sum q _ {i} + \sum m _ {i} + \sum p _ {i} \leq f} \lambda_ {(q _ {i}), (m _ {i}), (p _ {i})} u ((q _ {i}), (m _ {i}), (p _ {i})).
\]

Let \(v((q_i),(m_i),(p_i))\) be the monomial \(X_{-\alpha_1}^{q_1}\ldots X_{-\alpha_n}^{q_n}H_1^{m_1}\ldots H_l^{m_l}X_{\alpha_1}^{p_1}\ldots X_{\alpha_n}^{p_n}\) calculated in \(\mathbf{S}(\mathfrak{g})\) . Denoting by \(\mathbf{S}_d(\mathfrak{g})\) the sum of the homogeneous components of \(\mathbf{S}(\mathfrak{g})\) of degrees \(0,1,\ldots,d\) , we have

\[
\theta (z) \equiv \sum_ {\sum q _ {i} + \sum m _ {i} + \sum p _ {i} = f} \lambda_ {(q _ {i}), (m _ {i}), (p _ {i})} v ((q _ {i}), (m _ {i}), (p _ {i})) \pmod {\mathbf {S} _ {f - 1} (\mathfrak {g})}
\]

SO

\[
(\eta \circ \theta) (z) \equiv \sum_ {\sum m _ {i} = f} \lambda_ {(0), (m _ {i}), (0)} v ((0), (m _ {i}), (0)) \pmod {\mathbf {S} _ {f - 1} (\mathfrak {h})}
\]

and consequently

\[
(\eta \circ \theta) (z) \equiv \varphi (z) \pmod {\mathbf {S} _ {f - 1} (\mathfrak {h})}.\tag{3}
\]

\begin{enumerate}
\def\labelenumi{\alph{enumi})}
\setcounter{enumi}{2}
\tightlist
\item
  We show that \(\delta \circ \varphi : \mathbf{Z} \to \mathbf{S}(\mathfrak{h})^{\mathrm{W}}\) is bijective. The canonical filtrations on U and \(\mathbf{S}(\mathfrak{g})\) induce filtrations on Z, \(\mathrm{I}(\mathfrak{g})\) and \(\mathbf{S}(\mathfrak{h})^{\mathrm{W}}\) , and \(\theta, \eta\) are compatible with these filtrations, so that \(\mathrm{gr}(\eta\circ\theta)\) is an isomorphism from the vector space \(\mathrm{gr}(\mathbf{Z})\) to the vector space \(\mathrm{gr}(\mathbf{S}(\mathfrak{h})^{\mathrm{W}})\) . By (3), \(\mathrm{gr}(\varphi)=\mathrm{gr}(\eta\circ\theta)\) , and it is clear that \(\mathrm{gr}(\delta)\) is the identity. Hence \(\mathrm{gr}(\delta\circ\varphi)\) is bijective, so
\end{enumerate}

\[
\delta \circ \varphi : Z \to \mathbf {S} (\mathfrak {h}) ^ {W}
\]

is bijective (Commutative Algebra, Chap. III, §2, no. 8, Cor. 1 and 2 of Th. 1). d) Recall the notations in a). Let E be a simple g-module of highest weight \(\lambda\) , and \(\chi\) its central character (§6, no. 1, Def. 2). Let \(\varphi'\) and \(\delta'\) be the homomorphisms analogous to \(\varphi\) and \(\delta\) relative to the basis \(w(\mathrm{B})\) . The highest weight of E relative to \(w(\mathrm{B})\) is \(w(\lambda)\) . By §6, no. 4, Prop. 7,

\[
\varphi (u) (\lambda) = \chi (u) = \varphi^ {\prime} (u) (w \lambda)
\]

for all \(u\in \mathbf{Z}\) , so, by a),

\[
\begin{array}{c} (\delta \circ \varphi) (u) (w \lambda + w \rho) = (\delta \circ \varphi) (u) (\lambda + \rho) = \varphi (u) (\lambda) = \varphi^ {\prime} (u) (w \lambda) \\ = (\delta^ {\prime} \circ \varphi^ {\prime}) (u) (w \lambda + w \rho). \end{array}
\]

Thus, the polynomial functions \((\delta \circ \varphi)(u)\) and \((\delta' \circ \varphi')(u)\) coincide on \(w(\mathrm{P}_{++}) + w\rho\) , and hence are equal.

COROLLARY 1. For all \(\lambda \in \mathfrak{h}^*\) , let \(\chi_{\lambda}\) be the homomorphism \(z \mapsto (\varphi(z))(\lambda)\) from \(Z\) to \(k\) .

\begin{enumerate}
\def\labelenumi{(\roman{enumi})}
\item
  If \(k\) is algebraically closed, every homomorphism from \(Z\) to \(k\) is of the form \(\chi_{\lambda}\) for some \(\lambda \in \mathfrak{h}^*\) .
\item
  Let \(\lambda, \mu \in \mathfrak{h}^*\) . Then \(\chi_{\lambda} = \chi_{\mu}\) if and only if \(\mu + \rho \in \mathrm{W}(\lambda + \rho)\) .
\end{enumerate}

If \(k\) is algebraically closed, every homomorphism from \(\mathbf{S}(\mathfrak{h})^{\mathrm{W}}\) to \(k\) extends to a homomorphism from \(\mathbf{S}(\mathfrak{h})\) to \(k\) (Commutative Algebra, Chap. V, §1, no. 9, Prop. 22, and §2, no. 1, Cor. 4 of Th. 1), and every homomorphism from \(\mathbf{S}(\mathfrak{h})\) to \(k\) is of the form \(f \mapsto f(\lambda)\) for some \(\lambda \in \mathfrak{h}^*\) (Chap. VII, App. I, Prop. 1). Hence, if \(\chi\) is a homomorphism from \(Z\) to \(k\) , there exists (Th. 2) a \(\mu \in \mathfrak{h}^*\) such that, for all \(z \in Z\) ,

\[
\chi (z) = ((\delta \circ \varphi) (z)) (\mu) = (\varphi (z)) (\mu - \rho)
\]

hence (i).

Let \(\lambda, \mu \in \mathfrak{h}^*\) and assume that \(\chi_{\lambda} = \chi_{\mu}\) . Then, for all \(z \in \mathbf{Z}\) ,

\[
((\delta \circ \varphi) (z)) (\lambda + \rho) = (\varphi (z)) (\lambda) = \chi_ {\lambda} (z) = \chi_ {\mu} (z) = ((\delta \circ \varphi) (z)) (\mu + \rho);
\]

in other words, the homomorphisms from \(\mathbf{S}(\mathfrak{h})\) to k defined by \(\lambda + \rho\) and \(\mu + \rho\) coincide on \(\mathbf{S}(\mathfrak{h})^{\mathrm{W}}\) ; thus, assertion (ii) follows from Commutative Algebra, Chap. V, §2, no. 2, Cor. of Th. 2.

COROLLARY 2. Let E, E' be finite dimensional simple g-modules, and \(\chi, \chi'\) their central characters. If \(\chi = \chi'\) , E and E' are isomorphic.

Let \(\lambda, \lambda'\) be the highest weights of E, E'. By §6, no. 4, Prop. 7, \(\chi_{\lambda} = \chi = \chi' = \chi_{\lambda'}\) , so there exists \(w \in \mathrm{W}\) such that \(\lambda' + \rho = w(\lambda + \rho)\) . Since \(\lambda + \rho\) and \(\lambda' + \rho\) belong to the chamber defined by B, we have \(w = 1\) . Thus, \(\lambda = \lambda'\) , hence the corollary.

PROPOSITION 7. For any class \(\gamma\) of finite dimensional simple g-modules, let \(U_{\gamma}\) be the isotypical component of type \(\gamma\) of the g-module U (for the adjoint representation of g on U). Let \(\gamma_{0}\) be the class of the trivial g-module of dimension 1. Let \([U,U]\) be the vector subspace of U generated by the brackets of pairs of elements of U.

\begin{enumerate}
\def\labelenumi{(\roman{enumi})}
\item
  U is the direct sum of the \(\mathrm{U}_{\gamma}\) .
\item
  \(U_{\gamma_{0}} = Z\) , and \(\sum_{\gamma \neq \gamma_{0}} U_{\gamma} = [U, U]\) .
\item
  Let \(u \mapsto u^{\natural}\) be the projection of U onto Z defined by the decomposition \(\mathrm{U} = \mathrm{Z} \oplus [\mathrm{U}, \mathrm{U}]\) . If \(u \in \mathrm{U}\) and \(v \in \mathrm{U}\) , we have \((uv)^{\natural} = (vu)^{\natural}\) . If \(u \in \mathrm{U}\) and \(z \in \mathrm{Z}\) , we have \((uz)^{\natural} = u^{\natural}z\) .
\item
  Let \(\varphi\) be the Harish-Chandra homomorphism. Let \(\lambda \in \mathrm{P}_{++}\) , and let E be a finite dimensional simple \(\mathfrak{g}\) -module of highest weight \(\lambda\) . For all \(u \in \mathrm{U}\) , we have
\end{enumerate}

\[
\frac {1}{\dim \mathrm{E}} \mathrm{Tr} (u _ {\mathrm{E}}) = (\varphi (u ^ {\natural})) (\lambda).
\]

The \(\mathfrak{g}\) -module U is a direct sum of finite dimensional submodules. This implies (i).

It is clear that \(U_{\gamma_{0}} = Z\) . Let \(U'\) be a vector subspace of U defining a subrepresentation of class \(\gamma\) of the adjoint representation. Then either \([g, U'] = U'\) or \([g, U'] = 0\) . Thus, if \(\gamma \neq \gamma_{0}\) then \([g, U'] = U'\) , so \(\sum_{\gamma \neq \gamma_{0}} U_{\gamma} \subset [U, U]\) . On the other hand, if \(u \in U\) and \(x_{1}, \ldots, x_{n} \in g\) , then

\[
\begin{array}{c} [ x _ {1} \ldots x _ {n}, u ] = (x _ {1} \ldots x _ {n} u - x _ {2} \ldots x _ {n} u x _ {1}) + (x _ {2} \ldots x _ {n} u x _ {1} - x _ {3} \ldots x _ {n} u x _ {1} x _ {2}) \\ + \dots + (x _ {n} u x _ {1} \ldots x _ {n - 1} - u x _ {1} \ldots x _ {n}) \in [ \mathfrak {g}, \mathrm{U} ]. \end{array}
\]

Hence \([\mathrm{U},\mathrm{U}]\subset \left[\mathfrak{g},\sum_{\gamma}\mathrm{U}_{\gamma}\right] = \left[\mathfrak{g},\sum_{\gamma \neq \gamma_0}\mathrm{U}_\gamma \right]\subset \sum_{\gamma \neq \gamma_0}\mathrm{U}_\gamma .\) This proves (ii). Under these conditions, (iii) follows from Chap. I, §6, no. 9, Lemma 5.

Finally, let E, \(\lambda\) be as in (iv). Then

\[
\begin{array}{r l} & {\mathrm{Tr} (u _ {\mathrm{E}}) = \mathrm{Tr} ((u ^ {\natural}) _ {\mathrm{E}}) \qquad \text {since} u - u ^ {\natural} \in [ \mathrm{U}, \mathrm{U} ]} \\ & {\qquad = \mathrm{Tr} (\varphi (u ^ {\natural}) (\lambda). 1) \qquad (\S 6, \text {no.} 4, \text {Prop.} 7)} \\ & {\qquad = (\dim \mathrm{E}). \varphi (u ^ {\natural}) (\lambda).} \end{array}
\]

\section{THE FORMULA OF HERMANN WEYL}

In this paragraph, we retain the general notations of §6 and §7.

\subsection{CHARACTERS OF FINITE DIMENSIONAL g-MODULES}

Let \((e^{\lambda})_{\lambda \in \mathfrak{h}^{*}}\) be the canonical basis of the ring \(\mathbf{Z}[\mathfrak{h}^*]\) . Give the space \(\mathbf{Z}^{\mathfrak{h}^*}\) of all maps from \(\mathfrak{h}^*\) to \(\mathbf{Z}\) the product topology of the discrete topologies on the factors. If \(\varphi \in \mathbf{Z}^{\mathfrak{h}^*}\) , the family \((\varphi (\nu)e^{\nu})_{\nu \in \mathfrak{h}^*}\) is summable, and

\[
\varphi = \sum_ {\nu \in \mathfrak {h} ^ {*}} \varphi (\nu) e ^ {\nu}.
\]

Let \(\mathbf{Z}\langle\mathrm{P}\rangle\) be the set of \(\varphi\in\mathbf{Z}^{\mathfrak{h}^{*}}\) whose support is contained in a finite union of sets of the form \(\nu-\mathrm{P}_{+}\) , where \(\nu\in\mathfrak{h}^{*}\) . Then \(\mathbf{Z}[\mathrm{P}]\subset\mathbf{Z}\langle\mathrm{P}\rangle\subset\mathbf{Z}^{\mathfrak{h}^{*}}\) . Define on \(\mathbf{Z}\langle\mathrm{P}\rangle\) a ring structure extending that of \(\mathbf{Z}[\mathrm{P}]\) by putting, for \(\varphi,\psi\in\mathbf{Z}\langle\mathrm{P}\rangle\) and \(\nu\in\mathfrak{h}^{*}\) ,

\[
(\varphi \psi) (\nu) = \sum_ {\mu \in \mathfrak {h} ^ {*}} \varphi (\mu) \psi (\nu - \mu)
\]

(the family \((\varphi(\mu)\psi(\nu - \mu))_{\mu \in \mathfrak{h}^*}\) has finite support, in view of the condition satisfied by the supports of \(\varphi\) and \(\psi\) ). If \(\varphi = \sum_{\nu} x_{\nu} e^{\nu}\) and \(\psi = \sum_{\nu} y_{\nu} e^{\nu}\) , then \(\varphi \psi = \sum_{\nu, \mu} x_{\nu} y_{\mu} e^{\nu + \mu}\) .

Let \(\nu \in \mathfrak{h}^*\) . A partition of \(\nu\) into positive roots is a family \((n_{\alpha})_{\alpha \in \mathrm{R}_{+}}\) , where the \(n_{\alpha}\) are integers \(\geq 0\) such that \(\nu = \sum_{\alpha \in \mathrm{R}_{+}} n_{\alpha} \alpha\) . We denote by \(\mathfrak{P}(\nu)\) the number of partitions of \(\nu\) into positive roots. We have

\[
\mathfrak {P} (\nu) > 0 \Longleftrightarrow \nu \in \mathrm{Q} _ {+}.
\]

In this paragraph, we denote by K the following element of \(Z\langle P\rangle\) :

\[
\mathrm{K} = \sum_ {\gamma \in \mathrm{Q} _ {+}} \mathfrak {P} (\gamma) e ^ {- \gamma}.
\]

Now recall (Chap. VI, §3, no. 3, Prop. 2) that

\[
d = \prod_ {\alpha \in \mathrm{R} _ {+}} (e ^ {\alpha / 2} - e ^ {- \alpha / 2}) = \sum_ {w \in \mathrm{W}} \varepsilon (w) e ^ {w \rho}
\]

is an anti-invariant element of \(\mathbf{Z}[\mathrm{P}]\) .

Lemma 1. In the ring \(\mathbf{Z}\langle \mathrm{P}\rangle\) , we have \(\mathrm{K}\) . \(\prod_{\alpha \in \mathrm{R}_{+}}(1 - e^{-\alpha}) = \mathrm{Ke}^{-\rho}d = 1\) . Indeed,

\[
\mathrm{K} = \prod_ {\alpha \in \mathrm{R} _ {+}} \left(e ^ {0} + e ^ {- \alpha} + e ^ {- 2 \alpha} + \dots\right)
\]

SO

\[
K e ^ {- \rho} d = \prod_ {\alpha \in \mathrm{R} _ {+}} \left(1 + e ^ {- \alpha} + e ^ {- 2 \alpha} + \dots\right) \prod_ {\alpha \in \mathrm{R} _ {+}} \left(1 - e ^ {- \alpha}\right) = 1.
\]

Lemma 2. Let \(\lambda \in \mathfrak{h}^*\) . The module \(Z(\lambda)\) (§6, no. 3) admits a character that is an element of \(\mathbf{Z}\langle P\rangle\) , and we have \(d.\) ch \(Z(\lambda) = e^{\lambda + \rho}\) .

Let \(\alpha_{1},\ldots ,\alpha_{q}\) be distinct elements of \(\mathbb{R}_+\) . The \(X_{-\alpha_1}^{n_1}X_{-\alpha_2}^{n_2}\dots X_{-\alpha_q}^{n_q}\otimes 1\) form a basis of \(\mathrm{Z}(\lambda)\) (§6, Prop. 6 (iii)). For \(h\in \mathfrak{h}\) , we have

\[
\begin{array}{l} h. (X _ {- \alpha_ {1}} ^ {n _ {1}} X _ {- \alpha_ {2}} ^ {n _ {2}} \ldots X _ {- \alpha_ {q}} ^ {n _ {q}} \otimes 1) \\ \qquad = [ h, X _ {- \alpha_ {1}} ^ {n _ {1}} \ldots X _ {- \alpha_ {q}} ^ {n _ {q}} ] \otimes 1 + (X _ {- \alpha_ {1}} ^ {n _ {1}} \ldots X _ {- \alpha_ {q}} ^ {n _ {q}}) \otimes h. 1 \\ \qquad = (\lambda - n _ {1} \alpha_ {1} - \dots - n _ {q} \alpha_ {q}) (h) (X _ {- \alpha_ {1}} ^ {n _ {1}} \ldots X _ {- \alpha_ {q}} ^ {n _ {q}} \otimes 1). \end{array}
\]

Thus, the dimension of \(\mathrm{Z}(\lambda)^{\lambda - \mu}\) is \(\mathfrak{P}(\mu)\) . This proves that \(\operatorname{ch} \mathrm{Z}(\lambda)\) is defined, is an element of \(\mathbf{Z}\langle \mathrm{P} \rangle\) , and that

\[
\operatorname{ch} Z (\lambda) = \sum_ {\mu} \mathfrak {P} (\mu) e ^ {\lambda - \mu} = \mathrm{Ke} ^ {\lambda}.
\]

It now suffices to apply Lemma 1.

Lemma 3. Let M be a g-module which admits a character ch(M) whose support is contained in a finite union of the sets \(\mu-P_{+}\) . Let U be the enveloping algebra of g, Z the centre of U, \(\lambda_{0}\in h^{*}\) , and \(\chi_{\lambda_{0}}\) the corresponding homomorphism from Z to k ( \(\S8\) , Cor. 1 of Th. 2). Assume that, for all \(z\in Z\) , \(z_{M}\) is the homothety with ratio \(\chi_{\lambda_{0}}(z)\) . Let \(D_{M}\) be the set of \(\lambda\in W(\lambda_{0}+\rho)-\rho\) such that \(\lambda+Q_{+}\) meets Supp(ch M). Then ch(M) is a Z-linear combination of the ch \(Z(\lambda)\) for \(\lambda\in D_{M}\) .

If \(\mathrm{Supp(chM)}\) is empty, the lemma is clear. Assume that \(\mathrm{Supp(chM)} \neq \emptyset\) . Let \(\lambda\) be a maximal element of this support, and put \(\dim M^{\lambda} = m\) . There exists a \(\mathfrak{g}\) -homomorphism \(\varphi\) from \((Z(\lambda))^{m}\) to M which maps \((Z(\lambda)^{\lambda})^{m}\) bijectively onto \(M^{\lambda}\) ( \(\S 6\) , no. 3, Prop. 6 (i)). Thus, the central character of \(Z(\lambda)\) is \(\chi_{\lambda_0}\) , so \(\lambda \in W(\lambda_0 + \rho) - \rho\) ( \(\S 8\) , no. 5, Cor. 1 of Th. 2). This proves that \(D_M \neq \emptyset\) , and allows us to argue by induction on \(CardD_M\) . Let L and N be the kernel and cokernel of \(\varphi\) . Then we have an exact sequence of \(\mathfrak{g}\) -homomorphisms:

\[
0 \to \mathrm{L} \to (\mathrm{Z} (\lambda)) ^ {m} \to \mathrm{M} \to \mathrm{N} \to 0
\]

SO

\[
\operatorname{ch} (\mathrm{M}) = - \operatorname{ch} (\mathrm{L}) + m \operatorname{ch} \mathrm{Z} (\lambda) + \operatorname{ch} (\mathrm{N})
\]

(§7, no. 7, formula (6)). The sets Supp(ch L) and Supp(ch N) are contained in a finite union of sets \(\mu -\mathrm{P}_{+}\) . For \(z\in \mathbf{Z}\) , \(z_{\mathrm{L}}\) and \(z_{\mathrm{N}}\) are homotheties with ratio \(\chi_{\lambda_{0}}(z)\) . Clearly, \(D_{N} \subset D_{M}\) . On the other hand, \((\lambda + Q_{+}) \cap \text{Supp(ch M)} = \{\lambda\}\) , and \(\lambda \notin \text{Supp(ch N)}\) , so \(\lambda \notin D_{N}\) and

\[
\mathrm{Card} \mathrm{D} _ {\mathrm{N}} <   \mathrm{Card} \mathrm{D} _ {\mathrm{M}}.
\]

On the other hand, L is a submodule of \((Z(\lambda))^{m}\) ; if \(\lambda^{\prime} \in D_{L}\) , then \(\lambda^{\prime} + Q_{+}\) meets Supp(ch L) ⊂ Supp ch Z(λ), so \(\lambda \in \lambda^{\prime} + Q_{+}\) (§6, no. 1, Prop. 1 (ii)); it follows that \(D_{L} \subset D_{M}\) . Since \(L \cap (Z(\lambda)^{\lambda})^{m} = 0\) , we have \(\lambda \notin D_{L}\) , so

\(\mathrm{CardD_L} <   \mathrm{CardD_M}\)

It now suffices to apply the induction hypothesis.

THEOREM 1 (Character Formula of H. Weyl). Let M be a finite dimensional simple g-module, and \(\lambda\) its highest weight. Then

\[
\left(\sum_ {w \in \mathrm{W}} \varepsilon (w) e ^ {w \rho}\right). \mathrm{chM} = \sum_ {w \in \mathrm{W}} \varepsilon (w) e ^ {w (\lambda + \rho)}.
\]

With the notations of Lemma 3, the central character of M is \(\chi_{\lambda}\) ( \(\S6\) , no. 4, Prop. 7). Hence, by Lemmas 2 and 3, d.ch M is a Z-linear combination of the \(e^{\mu+\rho}\) such that

\[
\mu + \rho \in W (\lambda + \rho).
\]

On the other hand, by §7, no. 7, Lemma 7, d.ch M is anti-invariant, and its unique maximal term is \(e^{\lambda+\rho}\) , hence the theorem.

Example. Take \(\mathfrak{g} = \mathfrak{sl}(2,k)\) , \(\mathfrak{h} = kH\) . Let \(\alpha\) be the root of \((\mathfrak{g},\mathfrak{h})\) such that \(\alpha (H) = 2\) . The \(\mathfrak{g}\) -module \(V(m)\) has highest weight \((m / 2)\alpha\) . Hence

\[
\begin{array}{r l} & {\mathrm{ch} (\mathrm{V} (m)) = (e ^ {(m / 2) \alpha + \frac {1}{2} \alpha} - e ^ {- (m / 2) \alpha - \frac {1}{2} \alpha}) / (e ^ {\frac {1}{2} \alpha} - e ^ {- \frac {1}{2} \alpha})} \\ & {\qquad = e ^ {- (m / 2) \alpha}. (e ^ {(m + 1) \alpha} - 1) / (e ^ {\alpha} - 1)} \\ & {\qquad = e ^ {- (m / 2) \alpha} (e ^ {m \alpha} + e ^ {(m - 1) \alpha} + \dots + 1)} \\ & {\qquad = e ^ {(m / 2) \alpha} + e ^ {(m - 2) \alpha / 2} + \dots + e ^ {- (m / 2) \alpha}} \end{array}
\]

which also follows easily from §1, no. 2, Prop. 2.

\subsection{DIMENSIONS OF SIMPLE g-MODULES}

If \(\mu \in \mathfrak{h}^*\) , put \(\mathrm{J}(e^{\mu}) = \sum_{w\in \mathrm{W}}\varepsilon (w)e^{w\mu}\) , cf.~Chap. VI, §3, no. 3.

THEOREM 2. Let E be a finite dimensional simple g-module, \(\lambda\) its highest weight and \((\cdot|\cdot)\) a W-invariant non-degenerate positive symmetric bilinear form on \(h^{*}\) . Then:

\[
\dim \mathrm{E} = \prod_ {\alpha \in \mathrm{R} _ {+}} \frac {\langle \lambda + \rho , H _ {\alpha} \rangle}{\langle \rho , H _ {\alpha} \rangle} = \prod_ {\alpha \in \mathrm{R} _ {+}} \left(1 + \frac {(\lambda | \alpha)}{(\rho | \alpha)}\right).
\]

Let \(\mathrm{T}\) be an indeterminate. For all \(\nu \in \mathrm{P}\) , denote by \(f_{\nu}\) the homomorphism from \(\mathbf{Z}[\mathrm{P}]\) to \(\mathbf{R}[[\mathrm{T}]]\) that takes \(e^{\mu}\) to \(e^{(\nu |\mu)\mathrm{T}}\) for all \(\mu \in \mathrm{P}\) . Then \(\dim \mathrm{E}\) is the constant term of the series \(f_{\nu}(\operatorname{ch}\mathrm{E})\) .

For all \(\mu, \nu \in \mathrm{P}\) , we have

\[
\begin{array}{c} f _ {\nu} (\mathrm{J} (e ^ {\mu})) = \sum_ {w \in \mathrm{W}} \varepsilon (w) e ^ {(\nu | w \mu) \mathrm{T}} \\ = \sum_ {w \in \mathrm{W}} \varepsilon (w) e ^ {(w ^ {- 1} \nu | \mu) \mathrm{T}} = f _ {\mu} (\mathrm{J} (e ^ {\nu})). \end{array}
\]

In particular, in view of Chap. VI, §3, no. 3, formula (3),

\[
f _ {\rho} (\mathrm{J} (e ^ {\mu})) = f _ {\mu} (\mathrm{J} (e ^ {\rho})) = e ^ {(\mu | \rho) \mathrm{T}} \prod_ {\alpha \in \mathrm{R} _ {+}} (1 - e ^ {- (\mu | \alpha) \mathrm{T}}).
\]

Hence, setting \(\mathrm{Card}(\mathrm{R}_{+})=\mathrm{N}\) ,

\[
f _ {\rho} (\mathrm{J} (e ^ {\mu})) \equiv \mathrm{T} ^ {\mathrm{N}} \prod_ {\alpha \in \mathrm{R}} (\mu | \alpha) \pmod {\mathrm{T} ^ {\mathrm{N} + 1} \mathbf {R} [ [ \mathrm{T} ] ]}.
\]

The equality \(\mathrm{J}(e^{\lambda +\rho}) = \mathrm{ch}(\mathrm{E}).\mathrm{J}(e^{\rho})\) (Th. 1) thus implies that

\[
\mathrm{T} ^ {\mathrm{N}} \prod_ {\alpha \in \mathrm{R} _ {+}} (\lambda + \rho | \alpha) \equiv f _ {\rho} (\operatorname{ch} \mathrm{E}). \mathrm{T} ^ {\mathrm{N}} \prod_ {\alpha \in \mathrm{R} _ {+}} (\rho | \alpha) \pmod {\mathrm{T} ^ {\mathrm{N} + 1} \mathbf {R} [ [ \mathrm{T} ] ]}
\]

SO

\[
\dim \mathrm{E} = \left(\prod_ {\alpha \in \mathrm{R} _ {+}} (\lambda + \rho | \alpha)\right) / \left(\prod_ {\alpha \in \mathrm{R} _ {+}} (\rho | \alpha)\right) = \prod_ {\alpha \in \mathrm{R} _ {+}} \left(1 + \frac {(\lambda | \alpha)}{(\rho | \alpha)}\right).
\]

Now, if \(\alpha \in \mathrm{R}_{+}\) , \(\alpha\) can be identified with an element of \(\mathfrak{h}_{\mathbf{R}}\) proportional to \(H_{\alpha}\) , so

\[
(\lambda + \rho | \alpha) / (\rho | \alpha) = \langle \lambda + \rho , H _ {\alpha} \rangle / \langle \rho , H _ {\alpha} \rangle .
\]

Examples. 1) In the Example of no. 1, we find that

\[
\dim \mathrm{V} (m) = \left(\frac {m}{2} \alpha + \frac {\alpha}{2}\right) (H _ {\alpha}) / \frac {\alpha}{2} (H _ {\alpha}) = m + 1,
\]

which we knew in §1.

\begin{enumerate}
\def\labelenumi{\arabic{enumi})}
\setcounter{enumi}{1}
\tightlist
\item
  Take g to be the splittable simple Lie algebra of type \(G_{2}\) and adopt the notations of Chap. VI, Plate IX. Give \(h_{R}^{*}\) the W-invariant positive symmetric form \((\cdot|\cdot)\) such that \((\alpha_{1}|\alpha_{1})=1\) . Then \(\rho=\varpi_{1}+\varpi_{2}\) and
\end{enumerate}

\[
\begin{array}{c} (\varpi_ {1} | \alpha_ {1}) = \frac {1}{2}, (\varpi_ {1} | \alpha_ {2}) = 0, (\varpi_ {1} | \alpha_ {2} + \alpha_ {1}) = \frac {1}{2}, \\ (\varpi_ {1} | \alpha_ {2} + 2 \alpha_ {1}) = 1, (\varpi_ {1} | \alpha_ {2} + 3 \alpha_ {1}) = \frac {3}{2}, (\varpi_ {1} | 2 \alpha_ {2} + 3 \alpha_ {1}) = \frac {3}{2}, \\ (\varpi_ {2} | \alpha_ {1}) = 0, (\varpi_ {2} | \alpha_ {2}) = \frac {3}{2}, (\varpi_ {2} | \alpha_ {2} + \alpha_ {1}) = \frac {3}{2}, \\ (\varpi_ {2} | \alpha_ {2} + 2 \alpha_ {1}) = \frac {3}{2}, (\varpi_ {2} | \alpha_ {2} + 3 \alpha_ {1}) = \frac {3}{2}, (\varpi_ {2} | 2 \alpha_ {2} + 3 \alpha_ {1}) = 3. \end{array}
\]

Thus, if \(n_1, n_2\) are integers \(\geq 0\) , the dimension of the simple representation of highest weight \(n_1\varpi_1 + n_2\varpi_2\) is

\[
\begin{array}{c} \left(1 + \frac {n _ {1} / 2}{\frac {1}{2}}\right) \left(1 + \frac {3 n _ {2} / 2}{\frac {3}{2}}\right) \left(1 + \frac {n _ {1} / 2 + 3 n _ {2} / 2}{\frac {1}{2} + \frac {3}{2}}\right) \left(1 + \frac {n _ {1} + 3 n _ {2} / 2}{1 + \frac {3}{2}}\right) \\ \times \left(1 + \frac {3 n _ {1} / 2 + 3 n _ {2} / 2}{\frac {3}{2} + \frac {3}{2}}\right) \left(1 + \frac {3 n _ {1} / 2 + 3 n _ {2}}{\frac {3}{2} + 3}\right) \\ = (1 + n _ {1}) (1 + n _ {2}) \left(1 + \frac {n _ {1} + 3 n _ {2}}{4}\right) \left(1 + \frac {2 n _ {1} + 3 n _ {2}}{5}\right) \left(1 + \frac {n _ {1} + n _ {2}}{2}\right) \\ \times \left(1 + \frac {n _ {1} + 2 n _ {2}}{3}\right) \\ = \frac {(1 + n _ {1}) (1 + n _ {2}) (2 + n _ {1} + n _ {2}) (3 + n _ {1} + 2 n _ {2}) (4 + n _ {1} + 3 n _ {2}) (5 + 2 n _ {1} + 3 n _ {2})}{5 !}. \end{array}
\]

In particular, the fundamental representation of highest weight \(\varpi_{1}\) (resp. \(\varpi_{2}\) ) is of dimension 7 (resp. 14).

\subsection{MULTIPLICITIES OF WEIGHTS OF SIMPLE g-MODULES}

PROPOSITION 1. Let \(\omega \in \mathrm{P}_{++}\) . For all \(\lambda \in \mathrm{P}\) , the multiplicity of \(\lambda\) in \(\operatorname{E}(\omega)\) is

\[
m _ {\lambda} = \sum_ {w \in \mathrm{W}} \varepsilon (w) \mathfrak {P} (w (\omega + \rho) - (\lambda + \rho)).
\]

By Th. 1 and Lemma 1,

\[
\operatorname{ch} \mathrm{E} (\omega) = \mathrm{K} e ^ {- \rho} d \operatorname{ch} \mathrm{E} (\omega) = \mathrm{K} e ^ {- \rho} \sum_ {w \in \mathrm{W}} \varepsilon (w) e ^ {w (\omega + \rho)}
\]

SO

\[
\operatorname{ch} \mathrm{E} (\omega) = \sum_ {w \in \mathrm{W}, \gamma \in \mathrm{Q} _ {+}} \varepsilon (w) \mathfrak {P} (\gamma) e ^ {- \rho + w (\omega + \rho) - \gamma}
\]

and

\[
m _ {\lambda} = \sum_ {w \in \mathrm{W}, \gamma \in \mathrm{Q} _ {+}, \gamma = - \lambda - \rho + w (\omega + \rho)} \varepsilon (w) \mathfrak {P} (\gamma).
\]

COROLLARY. If \(\lambda\) is a weight of \(\operatorname{E}(\omega)\) distinct from \(\omega\) ,

\[
m _ {\lambda} = - \sum_ {w \in \mathrm{W}, w \neq 1} \varepsilon (w) m _ {\lambda + \rho - w \rho}.
\]

Apply Prop. 1 with \(\omega = 0\) . If \(\mu \in \mathrm{P} - \{0\}\) , we find that

\[
0 = \sum_ {w \in \mathrm{W}} \varepsilon (w) \mathfrak {P} (w \rho + \mu - \rho)
\]

hence

\[
\mathfrak {P} (\mu) = - \sum_ {w \in \mathrm{W}, w \neq 1} \varepsilon (w) \mathfrak {P} (\mu + w \rho - \rho).\tag{1}
\]

Prop. 1 also gives

\[
m _ {\lambda} = - \sum_ {w \in W} \varepsilon (w) \sum_ {w ^ {\prime} \in W, w ^ {\prime} \neq 1} \varepsilon (w ^ {\prime}) \mathfrak {P} (w (\omega + \rho) - (\lambda + \rho) + w ^ {\prime} \rho - \rho)
\]

since \(w(\omega + \rho) \neq \lambda + \rho\) for all \(w \in W\) (§7, Prop. 5 (iii)). Hence,

\[
\begin{array}{l} m _ {\lambda} = - \sum_ {w ^ {\prime} \in \mathrm{W}, w ^ {\prime} \neq 1} \varepsilon (w ^ {\prime}) \sum_ {w \in \mathrm{W}} \varepsilon (w) \mathfrak {P} (w (\omega + \rho) - (\lambda + \rho - w ^ {\prime} \rho + \rho)) \\ = - \sum_ {w ^ {\prime} \in \mathrm{W}, w ^ {\prime} \neq 1} \varepsilon (w ^ {\prime}) m _ {\lambda + \rho - w ^ {\prime} \rho} \quad (\text {Prop. 1}). \end{array}
\]

\subsection{DECOMPOSITION OF TENSOR PRODUCTS OF SIMPLE g-MODULES}

PROPOSITION 2. Let \(\lambda, \mu \in \mathrm{P}_{++}\) . In \(\mathcal{R}(\mathfrak{g})\) , we have

\[
[ \lambda ]. [ \mu ] = \sum_ {\nu \in \mathrm{P} _ {+ +}} m (\lambda , \mu , \nu) [ \nu ]
\]

with

\[
m (\lambda , \mu , \nu) = \sum_ {w, w ^ {\prime} \in \mathrm{W}} \varepsilon (w w ^ {\prime}) \mathfrak {P} (w (\lambda + \rho) + w ^ {\prime} (\mu + \rho) - (\nu + 2 \rho)).
\]

Let E, F be finite dimensional simple g-modules of highest weights \(\lambda, \mu\) . Let \(l_{\nu}\) be the length of the isotypical component of \(E \otimes F\) of highest weight \(\nu\) . It suffices to show that

\[
l _ {\nu} = \sum_ {w, w ^ {\prime} \in \mathrm{W}} \varepsilon (w w ^ {\prime}) \mathfrak {P} (w (\lambda + \rho) + w ^ {\prime} (\mu + \rho) - (\nu + 2 \rho)).\tag{2}
\]

Put \(c_{1} = \operatorname{ch}(\mathrm{E}) = \sum_{\sigma \in \mathrm{P}} m_{\sigma} e^{\sigma}\) , \(c_{2} = \operatorname{ch}(\mathrm{F})\) , and \(d = \mathrm{J}(e^{\rho})\) , where \(\mathrm{J}\) is defined as in no. 2. We have

\[
\sum_ {\xi \in \mathrm{P} _ {+ +}} l _ {\xi} \mathrm{ch} [ \xi ] = \mathrm{ch} (\mathrm{E} \otimes \mathrm{F}) = c _ {1} c _ {2}
\]

so, after multiplying by \(d\) and using Th. 1,

\[
\begin{array}{c} \sum_ {\xi \in \mathrm{P} _ {+ +}} l _ {\xi} \mathrm{J} (e ^ {\xi + \rho}) = c _ {1} \mathrm{J} (e ^ {\mu + \rho}) = \left(\sum_ {\sigma \in \mathrm{P}} m _ {\sigma} e ^ {\sigma}\right) \left(\sum_ {w \in \mathrm{W}} \varepsilon (w) e ^ {w (\mu + \rho)}\right) \\ = \sum_ {\tau \in \mathrm{P}} \left(\sum_ {w \in \mathrm{W}} \varepsilon (w) m _ {\tau + \rho - w (\mu + \rho)}\right) e ^ {\tau + \rho}. \end{array}\tag{3}
\]

Now, if \(\xi \in \mathrm{P}_{++}\) , \(\xi + \rho\) belongs to the chamber defined by B (Chap. VI, §1, no. 10); thus, for all \(w \in \mathrm{W}\) distinct from 1, we have \(w(\xi + \rho) \notin \mathrm{P}_{++}\) . Consequently, the coefficient of \(e^{\nu + \rho}\) in \(\sum_{\xi \in \mathrm{P}_{++}} l_{\xi} \mathrm{J}(e^{\xi + \rho})\) is equal to \(l_{\nu}\) . In view of (3), we obtain

\[
l _ {\nu} = \sum_ {w \in \mathrm{W}} \varepsilon (w) m _ {\nu + \rho - w (\mu + \rho)},
\]

that is, by Prop. 1,

\[
l _ {\nu} = \sum_ {w, w ^ {\prime} \in \mathrm{W}} \varepsilon (w) \varepsilon (w ^ {\prime}) \mathfrak {P} (w ^ {\prime} (\lambda + \rho) - (\nu + \rho - w (\mu + \rho) + \rho))
\]

which proves (2).

Example. We return to the Example of no. 1. Let \(\lambda = (n/2)\alpha, \mu = (p/2)\alpha\) , \(\nu = (q/2)\alpha\) with \(n \geq p\) . We have

\[
\begin{array}{l} m (\lambda , \mu , \nu) = \mathfrak {P} \left(\frac {n}{2} \alpha + \frac {\alpha}{2} + \frac {p}{2} \alpha + \frac {\alpha}{2} - \frac {q}{2} \alpha - \alpha\right) \\ \qquad - \mathfrak {P} \left(\frac {n}{2} \alpha + \frac {\alpha}{2} - \frac {p}{2} \alpha - \frac {\alpha}{2} - \frac {q}{2} \alpha - \alpha\right) \\ \qquad - \mathfrak {P} \left(- \frac {n}{2} \alpha - \frac {\alpha}{2} + \frac {p}{2} \alpha + \frac {\alpha}{2} - \frac {q}{2} \alpha - \alpha\right) \\ \qquad + \mathfrak {P} \left(- \frac {n}{2} \alpha - \frac {\alpha}{2} - \frac {p}{2} \alpha - \frac {\alpha}{2} - \frac {q}{2} \alpha - \alpha\right) \\ \qquad = \mathfrak {P} \left(\frac {n + p - q}{2} \alpha\right) - \mathfrak {P} \left(\frac {n - p - q - 2}{2} \alpha\right). \end{array}
\]

This is zero if \(n + p + q\) is not divisible by 2, or if \(q \geq n + p\) . If

\[
q = n + p - 2 r
\]

with \(r\) an integer \(\geq 0\) , we have

\[
m (\lambda , \mu , \nu) = \mathfrak {P} (r \alpha) - \mathfrak {P} ((r - p - 1) \alpha)
\]

hence \(m(\lambda,\mu,\nu)=1\) if \(r\leq p\) and \(m(\lambda,\mu,\nu)=0\) if r\textgreater p.~Finally, the g-module \(\mathrm{V}(n)\otimes\mathrm{V}(p)\) is isomorphic to

\[
\mathrm{V} (n + p) \oplus \mathrm{V} (n + p - 2) \oplus \mathrm{V} (n + p - 4) \oplus \dots \oplus \mathrm{V} (n - p)
\]

(Clebsch-Gordan formula).

\section{MAXIMAL SUBALGEBRAS OF SEMI-SIMPLE LIE ALGEBRAS}

THEOREM 1. Let V be a finite dimensional vector space, g a reductive Lie subalgebra in \(\mathfrak{gl}(V)\) , q a Lie subalgebra of g and \(\Phi\) the bilinear form \((x,y)\mapsto\mathrm{Tr}(xy)\) on \(g\times g\) . Assume that the orthogonal complement n of q with respect to \(\Phi\) is a Lie subalgebra of g consisting of nilpotent endomorphisms of V. Then q is a parabolic subalgebra of g.

\begin{enumerate}
\def\labelenumi{\alph{enumi})}
\tightlist
\item
  \(\mathfrak{q}\) is the normalizer of \(\mathfrak{n}\) in \(\mathfrak{g}\) : let \(\mathfrak{p}\) be this normalizer. Let \(x \in \mathfrak{q}\) and \(y \in \mathfrak{n}\) ; for all \(z \in \mathfrak{q}\) , we have \([z, x] \in \mathfrak{q}\) , so
\end{enumerate}

\[
\Phi ([ x, y ], z) = \Phi (y, [ z, x ]) = 0;
\]

in other words, \([x,y] \in n\) . Hence \(q \subset p\) . Since n is an ideal of p consisting of nilpotent endomorphisms of V, P is orthogonal to n with respect to \(\varPhi\) (Chap. I, no. 3, Prop. 4 d)). Since \(\varPhi\) is non-degenerate \(^{4}\) , \(p \subset q\) , hence our assertion.

\begin{enumerate}
\def\labelenumi{\alph{enumi})}
\setcounter{enumi}{1}
\tightlist
\item
  There exists a reductive Lie subalgebra m in \(\mathfrak{gl}(V)\) such that q is the semi-direct product of m and n: let \(\mathfrak{n}_{\mathrm{V}}(\mathfrak{q})\) be the largest ideal of q consisting of nilpotent endomorphisms of V. Then \(\mathfrak{n}_{\mathrm{V}}(\mathfrak{q})\) contains n, and it is orthogonal to q (loc. cit.); hence \(\mathfrak{n} = \mathfrak{n}_{\mathrm{V}}(\mathfrak{q})\) . Moreover, g is reductive in \(\mathfrak{gl}(V)\) by hypothesis, hence decomposable (Chap. VII, §5, no. 1, Prop. 2); since q is the intersection of g with the normalizer of n in \(\mathfrak{gl}(V)\) , it is a decomposable Lie algebra (loc. cit., Cor. 1 of Prop. 3). Thus, our assertion follows from Prop. 7 of Chap. VII, §5, no. 3.
\end{enumerate}

Choose a Cartan subalgebra \(\mathfrak{h}\) of \(\mathfrak{m}\) ; denote by \(\mathfrak{g}_1\) the commutant of \(\mathfrak{h}\) in \(\mathfrak{g}\) , and put \(\mathfrak{q}_1 = \mathfrak{q} \cap \mathfrak{g}_1, \mathfrak{n}_1 = \mathfrak{n} \cap \mathfrak{g}_1\) .

\begin{enumerate}
\def\labelenumi{\alph{enumi})}
\setcounter{enumi}{2}
\tightlist
\item
  The Lie algebras \(\mathfrak{g}_1, \mathfrak{q}_1\) and \(\mathfrak{n}_1\) satisfy the same hypotheses as \(\mathfrak{g}, \mathfrak{q}\) and \(\mathfrak{n}\) : since \(\mathfrak{m}\) is reductive in \(\mathfrak{gl}(\mathrm{V})\) , \(\mathfrak{h}\) is commutative and is composed of semi-simple endomorphisms of V (Chap. VII, §2, no. 4, Cor. 3 of Th. 2). Thus \(\mathfrak{g}_1 = \mathfrak{g}^0(\mathfrak{h})\) is reductive in \(\mathfrak{g}\) (Chap. VII, §1, no. 3, Prop. 11), hence also in \(\mathfrak{gl}(\mathrm{V})\) (Chap. I, §6, no. 6, Cor. 2 of Prop. 7). It is clear that \(\mathfrak{n}_1\) is composed of nilpotent endomorphisms of V. Since \(\mathfrak{h}\) is a subalgebra of \(\mathfrak{q}\) , reductive in \(\mathfrak{gl}(V)\) , the adjoint representation of \(\mathfrak{h}\) on \(\mathfrak{q}\) is semi-simple; by construction, \(\mathfrak{q}_1\) is the set of invariants of \(\mathrm{ad}_{\mathfrak{q}}(\mathfrak{h})\) , so \(\mathfrak{q} = \mathfrak{q}_1 + [\mathfrak{h}, \mathfrak{q}]\) (Chap. I, §3, no. 5, Prop. 6). Since
\end{enumerate}

\[
\Phi (\mathfrak {g} _ {1}, [ \mathfrak {h}, \mathfrak {q} ]) = \Phi ([ \mathfrak {h}, \mathfrak {g} _ {1} ], \mathfrak {q}) = 0,
\]

an element of \(\mathfrak{g}_1\) is orthogonal to \(\mathfrak{q}_1\) if and only if it is orthogonal to \(\mathfrak{q}\) ; consequently, \(\mathfrak{n}_1 = \mathfrak{g}_1 \cap \mathfrak{n}\) is the orthogonal complement of \(\mathfrak{q}_1\) in \(\mathfrak{g}_1\) .

\begin{enumerate}
\def\labelenumi{\alph{enumi})}
\setcounter{enumi}{3}
\item
  The Cartan subalgebra \(\mathfrak{h}\) of \(\mathfrak{m}\) is a Cartan subalgebra of \(\mathfrak{g}\) : We have \(\mathfrak{q} = \mathfrak{m} \oplus \mathfrak{n}\) and \(\mathfrak{h} = \mathfrak{m} \cap \mathfrak{g}_1\) , so it is immediate that \(\mathfrak{q}_1 = \mathfrak{h} \oplus \mathfrak{n}_1\) . Moreover, \([\mathfrak{h}, \mathfrak{n}_1] = 0\) , \(\mathfrak{h}\) is commutative and \(\mathfrak{n}_1\) is nilpotent, so the Lie algebra \(\mathfrak{q}_1\) is nilpotent. By \(a)\) and \(c)\) , \(\mathfrak{q}_1\) is the normalizer of \(\mathfrak{n}_1\) in \(\mathfrak{g}_1\) ; a fortiori, \(\mathfrak{q}_1\) is equal to its normalizer in \(\mathfrak{g}_1\) , hence is a Cartan subalgebra of \(\mathfrak{g}_1\) . Since \(\mathfrak{g}_1\) is reductive in \(\mathfrak{gl}(V)\) , it follows from Cor. 3 of Th. 2 of Chap. VII, §2, no. 4, that \(\mathfrak{q}_1\) is composed of semi-simple endomorphisms of V; thus, since \(\mathfrak{n}_1\) is composed of nilpotent endomorphisms of V, we have \(\mathfrak{n}_1 = 0\) . Consequently, \(\mathfrak{h} = \mathfrak{q}_1\) is a Cartan subalgebra of \(\mathfrak{g}_1\) , and since \(\mathfrak{g}_1\) normalizes \(\mathfrak{h}\) , we have \(\mathfrak{h} = \mathfrak{g}_1\) . Thus, we have proved that every element of \(\mathfrak{h}\) is a semi-simple element of \(\mathfrak{g}\) , and that the commutant of \(\mathfrak{h}\) in \(\mathfrak{g}\) is equal to \(\mathfrak{h}\) ; it follows that \(\mathfrak{h} = \mathfrak{g}^0(\mathfrak{h})\) , so \(\mathfrak{h}\) is a Cartan subalgebra of \(\mathfrak{g}\) .
\item
  \(\mathfrak{q}\) is a parabolic subalgebra of \(\mathfrak{g}\) : by the preceding, \(\mathfrak{h}\) is a Cartan subalgebra of \(\mathfrak{g}\) , \(\mathfrak{n}\) consists of nilpotent elements of \(\mathfrak{g}\) , and \([\mathfrak{h},\mathfrak{n}]\subset \mathfrak{n}\) . Let \(\bar{k}\) be an algebraic closure of \(k\) ; by definition, \(\mathfrak{q}\) is parabolic in \(\mathfrak{g}\) if and only if \(\bar{k}\otimes_k\mathfrak{q}\) is a parabolic subalgebra of \(\bar{k}\otimes_k\mathfrak{g}\) . The properties stated above being preserved by extension of scalars, for the proof we can restrict ourselves to the case in which \(\mathfrak{h}\) is splitting. Let R be the root system of \((\mathfrak{g},\mathfrak{h})\) ; by Prop. 2 (v) of §3, no. 1, there exists a subset P of R such that \(P\cap (-P) = \varnothing\) and \(\mathfrak{n} = \sum_{\alpha \in P}\mathfrak{g}^{\alpha}\) . Let \(P'\) be the set of roots \(\alpha\) such that \(-\alpha\notin P\) ; we have \(P' \cup (-P') = R\) , and the orthogonal complement \(\mathfrak{q}\) of \(\mathfrak{n}\) in \(\mathfrak{g}\) is equal to \(\mathfrak{h} + \sum_{\alpha \in P'}\mathfrak{g}^{\alpha}\) . We have proved that \(\mathfrak{q}\) is parabolic. Q.E.D.
\end{enumerate}

Lemma 1. Let g be a semi-simple Lie algebra, V a finite dimensional vector space, \(\rho\) a linear representation of g on V, D a vector subspace of V, h a Cartan subalgebra of g, s (resp. \(s'\) ) the set of \(x \in h\) such that \(\rho(x)\mathrm{D} \subset \mathrm{D}\) (resp. \(\rho(x)\mathrm{D} = 0\) ), and \(\Phi\) the bilinear form on g associated \(^{5}\) to \(\rho\) .

\begin{enumerate}
\def\labelenumi{(\roman{enumi})}
\item
  If \(\mathfrak{h}\) is splitting, the vector subspaces \(\mathfrak{s}\) and \(\mathfrak{s}'\) of \(\mathfrak{h}\) are rational over \(\mathbf{Q}\) .
\item
  If \(\rho\) is injective, the restriction of \(\Phi\) to \(\mathfrak{s}\) (resp. \(\mathfrak{s}'\) ) is non-degenerate.
\end{enumerate}

Assume that the Cartan subalgebra \(\mathfrak{h}\) is splitting. Let \(d\) be the dimension of D; put \(\mathrm{W} = \bigwedge^{d}(\mathrm{V})\) and \(\sigma = \bigwedge^{d}(\rho)\) ; denote also by \((e_1, \ldots, e_d)\) a basis of D and \(e = e_1 \wedge \cdots \wedge e_d\) a decomposable \(d\) -vector associated to D. Let P be the set of weights of \(\sigma\) with respect to \(\mathfrak{h}\) ; denote by \(\mathrm{W}^{\mu}\) the subspace of

W associated to the weight \(\mu\) , and put \(e = \sum_{\mu \in P} e^{\mu}\) (with \(e^{\mu} \in W^{\mu}\) for all \(\mu \in P\) ); finally, let \(P'\) be the set of weights \(\mu\) such that \(e^{\mu} \neq 0\) and let \(P''\) be the set of differences of elements of \(P'\) . Let x be in h; then x belongs to s if and only if there exists c in k such that \(\rho(x).e = c.e\) (Chap. VII, §5, no. 4, Lemma 2 (i)). Since \(\rho(x).e^{\mu} = \mu(x).e^{\mu}\) , we see that \(x \in s\) is equivalent to the relation `` \(\mu(x) = 0\) for all \(\mu \in P''\) ''. Now, the Q-structure of h is the Q-vector subspace \(h_{Q}\) of h generated by the coroots \(H_{\alpha}\) and all \(\mu\) in \(P''\) take rational values on \(h_{Q}\) ; it follows (Algebra, Chap. II, §8, no. 4, Prop. 5) that s is a subspace of h rational over Q.

For any weight \(\mu\in P\) , let \(p_{\mu}\) be the projection onto \(V^{\mu}\) associated to the decomposition \(V=\bigoplus_{\mu\in P}V^{\mu}\) ; denote by \(P_{1}\) the set of \(\mu\in P\) such that \(p_{\mu}(D)\neq0\) . It is immediate that \(s'\) is the intersection of the kernels (in h) of the elements of \(P_{1}\) ; it follows, in the same way as for s, that \(s'\) is a subspace of h rational over Q. This proves (i).

By extension of scalars, it suffices to prove (ii) when \(k\) is algebraically closed, hence when \(\mathfrak{h}\) is splitting. Let \(\mathfrak{m}\) be a vector subspace of \(\mathfrak{h}\) rational over \(\mathbf{Q}\) ; for all non-zero \(x\) in \(\mathfrak{m}_{\mathbf{Q}} = \mathfrak{m} \cap \mathfrak{h}_{\mathbf{Q}}\) , we have \(\varPhi(x, x) > 0\) by the Cor. of Prop. 1 of §7, no. 1. The restriction of \(\varPhi\) to \(\mathfrak{m}_{\mathbf{Q}}\) is non-degenerate, and hence so is the restriction of \(\varPhi\) to \(\mathfrak{m}\) since \(\mathfrak{m}\) is canonically isomorphic to \(k \otimes_{\mathbf{Q}} \mathfrak{m}_{\mathbf{Q}}\) .

DEFINITION 1. Let q be a Lie subalgebra of the semi-simple Lie algebra g. Then q is said to be decomposable in g if, for all \(x \in q\) , the semi-simple and nilpotent components of x in g belong to q. Denote by \(\mathfrak{n}_{\mathfrak{g}}(\mathfrak{q})\) the set of elements x of the radical of q such that \(ad_{g}x\) is nilpotent.

Let \(\rho\) be an injective representation of \(\mathfrak{g}\) on a finite dimensional vector space V. We know (Chap. I, §6, no. 3, Th. 3) that an element \(x\) of \(\mathfrak{g}\) is semi-simple (resp. nilpotent) if and only if the endomorphism \(\rho(x)\) of V is semi-simple (resp. nilpotent). It follows immediately that the algebra \(\mathfrak{q}\) is decomposable in \(\mathfrak{g}\) if and only if \(\rho(\mathfrak{q})\) is a decomposable subalgebra of \(\mathfrak{gl}(V)\) in the sense of Definition 1 of Chap. VII, §5, no. 1. With the notations of Chap. VII, §5, no. 3, we also have

\[
\rho (\mathfrak {n} _ {\mathfrak {g}} (\mathfrak {q})) = \mathfrak {n} _ {\mathrm{V}} (\rho (\mathfrak {q})).
\]

THEOREM 2. Let g be a semi-simple Lie algebra, n a subalgebra of g consisting of nilpotent elements, q the normalizer of n in g. Assume that n is the set of nilpotent elements of the radical of q. Then q is parabolic.

Note first of all that q is decomposable (Chap. VII, §5, no. 1, Cor. 1 of Prop. 3). By Th. 1, it suffices to prove that q is the orthogonal complement \(n^{0}\) of n with respect to the Killing form \(\varPhi\) of g. We know that \(q \subset n^{0}\) (Chap. I, §4, no. 3, Prop. 4 d)). By Chap. VII, §5, no. 3, Prop. 7, there exists a subalgebra m of q, reductive in g, such that q is the semi-direct product of m and n.~We show that the restriction of \(\varPhi\) to m is non-degenerate. Let c be the centre of m. We have \(\Phi([\mathfrak{m},\mathfrak{m}],\mathfrak{c})=0\) by Chap. I, §5, no. 5, Prop. 5, and the restriction of \(\Phi\) to \([m,m]\) is non-degenerate by Chap. I, §6, no. 1, Prop. 1. It remains to see that the restriction of \(\Phi\) to c is non-degenerate. Let k be a Cartan subalgebra of \([m,m]\) ; then \(k\oplus c\) is commutative and reductive in g. Let h be a Cartan subalgebra of g containing \(k\oplus c\) (Chap. VII, §2, no. 3, Prop. 10). Then \(h\cap q\) is a commutative subalgebra of q containing \(k\oplus c\) , and \(ad_{q}x\) is semi-simple for all \(x\in h\cap q\) ; hence \(h\cap q\) is contained in a Cartan subalgebra \(h'\) of q (Chap. VII, §2, no. 3, Prop. 10); let f be the projection of q onto m with kernel n; then \(f(h')\) is a Cartan subalgebra of m (Chap. VII, §2, no. 1, Cor. 2 of Prop. 4) containing \(k\oplus c\) , and consequently equal to \(k\oplus c\) ; this proves that \(f(h\cap q)=k\oplus c\) , and since every element of h is semi-simple in g, we have \(h\cap q=k\oplus c\) . Thus,

\[
\mathfrak {c} = \{x \in \mathfrak {h} \mid [ x, \mathfrak {n} ] \subset \mathfrak {n} \text {   and   } [ x, [ \mathfrak {m}, \mathfrak {m} ] ] = 0 \}.
\]

By Lemma 1, the restriction of \(\varPhi\) to \(\mathfrak{c}\) is non-degenerate.

Let \(\mathfrak{q}^0\) be the orthogonal complement of \(\mathfrak{q}\) in \(\mathfrak{g}\) relative to \(\varPhi\) . The preceding proves that \(\mathfrak{q} \cap \mathfrak{q}^0 = \mathfrak{n}\) . Assume that \(\mathfrak{q} \neq \mathfrak{q}^0\) , so \(\mathfrak{q}^0 \neq \mathfrak{n}\) (and \(\mathfrak{q}^0 \supset \mathfrak{n}\) ). Since \(\operatorname{ad}_{\mathfrak{g}} \mathfrak{n}\) leaves \(\mathfrak{q}\) stable, \(\operatorname{ad}_{\mathfrak{g}} \mathfrak{n}\) leaves \(\mathfrak{q}^0\) stable; Engel's theorem proves that there exists \(x \in \mathfrak{q}^0\) such that \(x \notin \mathfrak{n}\) and \([x, \mathfrak{n}] \subset \mathfrak{n}\) . But then \(x \in \mathfrak{q}^0 \cap \mathfrak{q} = \mathfrak{n}\) , a contradiction. Hence \(\mathfrak{q} = \mathfrak{n}^0\) .

COROLLARY 1. Let \(\mathfrak{q}\) be a maximal element of the set of subalgebras of \(\mathfrak{g}\) distinct from \(\mathfrak{g}\) . Then \(\mathfrak{q}\) is either parabolic or reductive in \(\mathfrak{g}\) .

We can assume that g is a Lie subalgebra of \(\mathfrak{gl}(V)\) for some finite dimensional vector space V. Let \(\mathfrak{e}(\mathfrak{q}) \subset \mathfrak{g}\) be the decomposable envelope of q. If \(\mathfrak{e}(\mathfrak{q}) = \mathfrak{g}\) , q is an ideal of g (Chap. VII, §5, no. 2, Prop. 4), hence is semi-simple, and consequently q is reductive in g. Assume that \(\mathfrak{e}(\mathfrak{q}) \neq \mathfrak{g}\) . Then \(\mathfrak{e}(\mathfrak{q}) = \mathfrak{q}\) , so q is decomposable. Assume that q is not reductive in g. Let n be the set of nilpotent elements of the radical of q. Then \(n \neq 0\) (Chap. VII, §5, no. 3, Prop. 7 (i)). Let p be the normalizer of n in g. Then \(p \supset q\) , and \(p \neq g\) since g is semi-simple. Hence p = q. Thus q is parabolic (Th. 1).

COROLLARY 2. Let \(\mathfrak{n}\) be a subalgebra of \(\mathfrak{g}\) consisting of nilpotent elements. There exists a parabolic subalgebra \(\mathfrak{q}\) of \(\mathfrak{g}\) with the following properties:

\begin{enumerate}
\def\labelenumi{(\roman{enumi})}
\item
  \(\mathfrak{n}\subset \mathfrak{n}_{\mathfrak{g}}(\mathfrak{q});\)
\item
  the normalizer of \(\mathfrak{n}\) in \(\mathfrak{g}\) is contained in \(\mathfrak{q}\) ;
\item
  every automorphism of g leaving n invariant leaves q invariant.
\end{enumerate}

If \(\mathfrak{g}\) is splittable, \(\mathfrak{n}\) is contained in a Borel subalgebra of \(\mathfrak{g}\) .

Let \(\mathfrak{q}_1\) be the normalizer of \(\mathfrak{n}\) in \(\mathfrak{g}\) . This is a decomposable subalgebra of \(\mathfrak{g}\) . Let \(\mathfrak{n}_1 = \mathfrak{n}_{\mathfrak{g}}(\mathfrak{q}_1)\) . Define inductively \(\mathfrak{q}_i\) to be the normalizer of \(\mathfrak{n}_{i-1}\) in \(\mathfrak{g}\) , and \(\mathfrak{n}_i\) to be equal to \(\mathfrak{n}_{\mathfrak{g}}(\mathfrak{q}_i)\) . The sequences \((\mathfrak{n}, \mathfrak{n}_1, \mathfrak{n}_2, \ldots)\) and \((\mathfrak{q}_1, \mathfrak{q}_2, \ldots)\) are increasing. There exists \(j\) such that \(\mathfrak{q}_j = \mathfrak{q}_{j+1}\) , in other words \(\mathfrak{q}_j\) is the normalizer of \(\mathfrak{n}_{\mathfrak{g}}(\mathfrak{q}_j)\) in \(\mathfrak{g}\) . Thus \(\mathfrak{q}_j\) is parabolic (Th. 1). We have \(\mathfrak{n} \subset \mathfrak{n}_j =\) \(\mathfrak{n}_{\mathfrak{g}}(\mathfrak{q}_{j})\) , and \(q_{1} \subset q_{j}\) ; every automorphism of g leaving n invariant evidently leaves \(n_{1}, n_{2}, \ldots\) and \(q_{1}, q_{2}, \ldots\) invariant. If g is splittable, \(q_{j}\) contains a Borel subalgebra b, and consequently (§3, no. 4, Prop. 13), we have \(\mathfrak{b} \supset \mathfrak{n}_{\mathfrak{g}}(\mathfrak{q}_{j}) \supset \mathfrak{n}\) .

THEOREM 3. Assume that k is algebraically closed. Let g be a semi-simple Lie algebra. Let a be a solvable subalgebra of g. There exists a Borel subalgebra of g containing a.

By Chap. VII, §5, no. 2, Cor. 1 (ii) of Prop. 4, we can assume that \(\mathfrak{a}\) is decomposable. There exists a commutative subalgebra \(\mathfrak{t}\) of \(\mathfrak{g}\) , consisting of semi-simple elements, such that \(\mathfrak{a}\) is the semi-direct product of \(\mathfrak{t}\) and \(\mathfrak{n}_{\mathfrak{g}}(\mathfrak{a})\) (Chap. VII, §5, no. 3, Cor. 2 of Prop. 6). There exists (Cor. 2 of Th. 2) a parabolic subalgebra \(\mathfrak{q}\) of \(\mathfrak{g}\) such that \(\mathfrak{n}_{\mathfrak{g}}(\mathfrak{a}) \subset \mathfrak{n}_{\mathfrak{g}}(\mathfrak{q})\) , and such that the normalizer of \(\mathfrak{n}_{\mathfrak{g}}(\mathfrak{a})\) in \(\mathfrak{g}\) is contained in \(\mathfrak{q}\) ; a fortiori, \(\mathfrak{a} \subset \mathfrak{q}\) . Let \(\mathfrak{b}\) be a Borel subalgebra of \(\mathfrak{g}\) contained in \(\mathfrak{q}\) and \(\mathfrak{h}\) a Cartan subalgebra of \(\mathfrak{g}\) contained in \(\mathfrak{b}\) . Then \(\mathfrak{h}\) is a Cartan subalgebra of \(\mathfrak{q}\) , so there exists \(s \in \mathrm{Aut}_e(\mathfrak{q})\) such that \(s(\mathfrak{t}) \subset \mathfrak{h}\) (Chap. VII, §2, no. 3, Prop. 10 and Chap. VII, §3, no. 2, Th. 1). We have \(s(\mathfrak{n}_{\mathfrak{g}}(\mathfrak{q})) = \mathfrak{n}_{\mathfrak{g}}(\mathfrak{q})\) (Chap. VII, §3, no. 1, Remark 1), so

\[
s (\mathfrak {a}) = s (t) + s \left(\mathfrak {n} _ {\mathfrak {g}} (\mathfrak {a})\right) \subset \mathfrak {h} + s \left(\mathfrak {n} _ {\mathfrak {g}} (\mathfrak {q})\right) = \mathfrak {h} + \mathfrak {n} _ {\mathfrak {g}} (\mathfrak {q}) \subset \mathfrak {b}.
\]

COROLLARY. If \(k\) is algebraically closed, every maximal solvable subalgebra of \(\mathfrak{g}\) is a Borel subalgebra.

\section{\texorpdfstring{CLASSES OF NILPOTENT ELEMENTS AND \(\mathfrak{sl}_2\) -TRIPLETS}{CLASSES OF NILPOTENT ELEMENTS AND \textbackslash mathfrak\{sl\}\_2 -TRIPLETS}}

In this paragraph, g denotes a finite dimensional Lie algebra.

\subsection{\texorpdfstring{DEFINITION OF \(\mathfrak{sl}_2\) -TRIPLETS}{DEFINITION OF \textbackslash mathfrak\{sl\}\_2 -TRIPLETS}}

DEFINITION 1. An \(\mathfrak{sl}_2\) -triplet in \(\mathfrak{g}\) is a sequence \((x, h, y)\) of elements of \(\mathfrak{g}\) , distinct from \((0, 0, 0)\) , such that

\[
[ h, x ] = 2 x, [ h, y ] = - 2 y, [ x, y ] = - h.
\]

Let \((x,h,y)\) be an \(\mathfrak{sl}_{2}\) -triplet in g. The linear map \(\tau\) from \(\mathfrak{sl}(2,k)\) to g such that \(\tau(X_{+}) = x, \tau(H) = h, \tau(X_{-}) = y\) is a homomorphism which is non-zero and hence injective (since \(\mathfrak{sl}(2,k)\) is simple), and with image \(kx + kh + ky\) . We thus obtain a canonical bijection from the set of \(\mathfrak{sl}_{2}\) -triplets in g to the set of injective homomorphisms from \(\mathfrak{sl}(2,k)\) to g. If g is semi-simple and if \((x,h,y)\) is an \(\mathfrak{sl}_{2}\) -triplet in g, then x and y are nilpotent elements of g and h is a semi-simple element of g (Chap. I, §6, no. 3, Prop. 4).

Lemma 1. Let \(x, h, y, y' \in \mathfrak{g}\) . If \((x, h, y)\) and \((x, h, y')\) are \(\mathfrak{sl}_2\) -triplets in \(\mathfrak{g}\) , then \(y = y'\) .

Indeed, \(y - y' \in \mathrm{Ker}(\mathrm{ad}_{\mathfrak{g}}x)\) and \((\mathrm{ad}_{\mathfrak{g}}h)(y - y') = -2(y - y')\) . But \(\mathrm{ad}_{\mathfrak{g}}x\) is injective on \(\mathrm{Ker}(p + \mathrm{ad}_{\mathfrak{g}}h)\) for every integer \(p > 0\) (§1, no. 2, Cor. of Prop. 2).

Lemma 2. Let \(\mathfrak{n}\) be a subalgebra of \(\mathfrak{g}\) such that, for all \(n\in \mathfrak{n}\) , \(\operatorname{ad}_{\mathfrak{g}}(n)\) is nilpotent. Let \(h\in \mathfrak{g}\) be such that \([h,\mathfrak{n}] = \mathfrak{n}\) . Then \(e^{\operatorname{ad}_{\mathfrak{g}}\mathfrak{n}}.h = h + \mathfrak{n}\) .

It is clear that \(e^{\operatorname{ad}_{\mathfrak{g}}(\mathfrak{n})}.h \subset h + \mathfrak{n}\) . We shall prove that, if \(v \in n\) , then \(h + v \in e^{\operatorname{ad}_{\mathfrak{g}}(\mathfrak{n})}.h\) . It suffices to prove that \(h + v \in e^{\operatorname{ad}_{\mathfrak{g}}(\mathfrak{n})}.h + \mathscr {C} ^ {p} \mathfrak {n}\) for all \(p \geq 1\) (since \(C^{p}n = 0\) for sufficiently large p). This is clear for p = 1 since \(C^{1}n = n\) . Assume now that we have proved the existence of \(y_{p} \in n\) and \(z_{p} \in C^{p}n\) such that \(h + v = e^{\operatorname{ad}_{\mathfrak{g}}y_{p}}.h + z_{p}\) . Since \((\operatorname{ad}_{\mathfrak{g}}h)(\mathfrak{n}) = \mathfrak{n}\) , \((\operatorname{ad}_{\mathfrak{g}}h)|\mathfrak{n}\) is a bijection from n to n, hence its restriction to \(C^{p}n\) , which leaves \(C^{p}n\) stable, is also bijective; consequently, there exists \(z \in C^{p}n\) such that \(z_{p} = [z, h]\) . Then

\[
e ^ {\operatorname{ad} _ {\mathfrak {g}} \left(y _ {p} + z\right)} h - e ^ {\operatorname{ad} _ {\mathfrak {g}} y _ {p}} h \in [ z, h ] + \mathscr {C} ^ {p + 1} \mathfrak {n}
\]

SO

\[
e ^ {\operatorname{ad} _ {\mathfrak {g}} \left(y _ {p} + z\right)} h \in h + v - z _ {p} + [ z, h ] + \mathscr {C} ^ {p + 1} \mathfrak {n} = h + v + \mathscr {C} ^ {p + 1} \mathfrak {n}
\]

which establishes our assertion by induction on p.

Lemma 3. Let \(x \in \mathfrak{g}\) , \(\mathfrak{p} = \operatorname{Ker}(\operatorname{ad}x)\) , \(\mathfrak{q} = \operatorname{Im}(\operatorname{ad}x)\) . Then \([\mathfrak{p}, \mathfrak{q}] \subset \mathfrak{q}\) , and \(\mathfrak{p} \cap \mathfrak{q}\) is a subalgebra of \(\mathfrak{g}\) .

If \(u \in \mathfrak{p}\) and \(v \in \mathfrak{q}\) , there exists \(w \in \mathfrak{g}\) such that \(v = [x, w]\) , so

\[
[ u, v ] = [ u, [ x, w ] ] = [ x, [ u, w ] ] - [ [ x, u ], w ] = [ x, [ u, w ] ] \in \mathfrak {q}.
\]

On the other hand, p is a subalgebra of g, so \([p \cap q, p \cap q] \subset p \cap q\) .

Lemma 4. Let \((x,h,y)\) and \((x,h',y')\) be \(\mathfrak{sl}_2\) -triplets in \(\mathfrak{g}\) . There exists \(z\in \mathfrak{g}\) such that \(\mathrm{ad}_{\mathfrak{g}}z\) is nilpotent and such that

\[
e ^ {\mathrm{ad} _ {\mathfrak {g}} z} x = x, \quad e ^ {\mathrm{ad} _ {\mathfrak {g}} z} h = h ^ {\prime}, \quad e ^ {\mathrm{ad} _ {\mathfrak {g}} z} y = y ^ {\prime}.
\]

Let \(\mathfrak{n} = \mathrm{Ker}(\mathrm{ad} x) \cap \mathrm{Im}(\mathrm{ad} x)\) . For all \(p \in \mathbf{Z}\) , let \(\mathfrak{g}_p = \mathrm{Ker}(\mathrm{ad} h - p)\) . By §1, no. 3 (applied to the adjoint representation of \(kx + ky + kh\) on \(\mathfrak{g}\) ), we have that \(\mathfrak{n} \subset \sum_{p > 0} \mathfrak{g}_p\) , so \(\mathrm{ad}_{\mathfrak{g}} n\) is nilpotent for all \(n \in \mathfrak{n}\) , and \([h, \mathfrak{n}] = \mathfrak{n}\) . We have \([x, h' - h] = 0\) and \([x, y - y'] = h' - h\) , so \(h' - h \in \mathfrak{n}\) . By Lemmas 2 and 3, there exists \(z \in \mathfrak{n}\) such that \(e^{\mathrm{ad}_{\mathfrak{g}} z} h = h'\) . Since \(z \in \mathrm{Ker} \mathrm{ad}_{\mathfrak{g}} x\) , we have \(e^{\mathrm{ad}_{\mathfrak{g}} z} x = x\) . Lemma 1 now proves that \(e^{\mathrm{ad}_{\mathfrak{g}} z} y = y'\) . Q.E.D.

Let G be a group of automorphisms of g. Then two \(sl_{2}\) -triplets \((x,h,y)\) , \((x',h',y')\) are said to be G-conjugate if there exists \(g\in G\) such that \(gx=x'\) , \(gh=h',gy=y'\) .

PROPOSITION 1. Let G be a group of automorphisms of g containing \(\mathrm{Aut}_e(\mathfrak{g})\) . Let \((x,h,y)\) and \((x',h',y')\) be \(\mathfrak{sl}_2\) -triplets in g. Let

\[
\mathfrak {t} = k x + k h + k y, \quad \mathfrak {t} ^ {\prime} = k x ^ {\prime} + k h ^ {\prime} + k y ^ {\prime}.
\]

Consider the following conditions:

\begin{enumerate}
\def\labelenumi{(\roman{enumi})}
\item
  x and \(x'\) are G-conjugate;
\item
  \((x, h, y)\) and \((x', h', y')\) are G-conjugate;
\item
  \(\mathfrak{t}\) and \(\mathfrak{t}'\) are G-conjugate.
\end{enumerate}

We have (i) \(\Longleftrightarrow\) (ii) \(\Longrightarrow\) (iii). If \(k\) is algebraically closed, the three conditions are equivalent.

\begin{enumerate}
\def\labelenumi{(\roman{enumi})}
\item
  \(\Longleftrightarrow\) (ii): This follows from Lemma 4.
\item
  \(\Longrightarrow\) (iii): This is clear.
\end{enumerate}

We assume that \(k\) is algebraically closed and prove that (iii) \(\Longrightarrow\) (i). We treat first the case in which \(\mathfrak{t} = \mathfrak{t}' = \mathfrak{g} = \mathfrak{sl}(2, k)\) . Since \(\operatorname{ad}_{\mathfrak{g}} x\) is nilpotent, the endomorphism \(x\) of \(k^2\) is nilpotent (Chap. I, §6, Th. 3), so there exists a matrix \(A \in \mathbf{GL}(2, k)\) such that \(AxA^{-1} = X\) , and consequently an automorphism \(\alpha\) of \(\mathfrak{sl}(2, k)\) such that \(\alpha(x) = x'\) ; now \(\alpha \in \operatorname{Aut}_e(\mathfrak{g})\) (§5, no. 3, Cor. 2 of Prop. 5). We now pass to the general case; we assume that \(\mathfrak{t}\) and \(\mathfrak{t}'\) are G-conjugate and prove that \(x\) and \(x'\) are G-conjugate. We can assume that \(\mathfrak{t} = \mathfrak{t}'\) . By the preceding, there exists \(\beta \in \operatorname{Aut}_e(\mathfrak{t})\) such that \(\beta x = x'\) . Now, if \(t \in \mathfrak{t}\) is such that \(\operatorname{ad}_{\mathfrak{t}} t\) is nilpotent, then \(\operatorname{ad}_{\mathfrak{g}} t\) is nilpotent; so \(\beta\) extends to an element of \(\operatorname{Aut}_e(\mathfrak{g})\) .

Remark. The three conditions of Prop. 1 are equivalent if we assume only that \(k = k^2\) (cf.~Exerc. 1).

\subsection{\texorpdfstring{\(\mathfrak{sl}_2\) -TRIPLETS IN SEMI-SIMPLE LIE ALGEBRAS}{\textbackslash mathfrak\{sl\}\_2 -TRIPLETS IN SEMI-SIMPLE LIE ALGEBRAS}}

Lemma 5. Let V be a finite dimensional vector space, A and B endomorphisms of V. Assume that A is nilpotent and that \([A, [A, B]] = 0\) . Then AB is nilpotent.

Put C = {[}A, B{]}. Since {[}A, C{]} = 0,

\[
[ \mathrm{A}, \mathrm{BC} ^ {p} ] = [ \mathrm{A}, \mathrm{B} ] \mathrm{C} ^ {p} = \mathrm{C} ^ {p + 1}
\]

for every integer \(p \geq 0\) . Consequently, \(\mathrm{Tr}(\mathrm{C}^{p}) = 0\) for \(p \geq 1\) , which proves that C is nilpotent (Algebra, Chap. VII, §3, no. 5, Cor. 4 of Prop. 13). Now let \(\bar{k}\) be an algebraic closure of k, and let \(\lambda \in \bar{k}\) , \(x \in V \otimes_{k} \bar{k}\) be such that ABx = \(\lambda x\) , \(x \neq 0\) . The relation \([[B, A], A] = 0\) shows that \([B, A^{p}] = p[B, A]A^{p-1}\) for every integer \(p \geq 0\) . Let r be the smallest integer such that \(A^{r}x = 0\) . Then

\[
\lambda \mathrm{A} ^ {r - 1} x = \mathrm{A} ^ {r - 1} \mathrm{AB} x = \mathrm{A} ^ {r} \mathrm{B} x = \mathrm{BA} ^ {r} x - [ \mathrm{B}, \mathrm{A} ^ {r} ] x = - r [ \mathrm{B}, \mathrm{A} ] \mathrm{A} ^ {r - 1} x.
\]

Since {[}B, A{]} is nilpotent and since \(A^{r-1}x \neq 0\) , this proves that \(\lambda = 0\) . Thus, all the eigenvalues of AB are zero, hence the lemma.

Lemma 6. Let \(h, x \in \mathfrak{g}\) be such that \([h, x] = 2x\) and \(h \in (\operatorname{ad} x)(\mathfrak{g})\) . Then there exists \(y \in \mathfrak{g}\) such that \((x, h, y)\) is either \((0, 0, 0)\) or an \(\mathfrak{sl}_2\) -triplet.

Let \(\mathfrak{g}'\) be the solvable Lie algebra \(kh + kx\) . Since \(x \in [\mathfrak{g}', \mathfrak{g}']\) , \(\operatorname{ad}_{\mathfrak{g}} x\) is nilpotent (Chap. I, §5, no. 3, Th. 1); let \(\mathfrak{n}\) be its kernel. Since {[}ad \(h\) , ad \(x] = 2\) ad \(x\) , we have (ad \(h\) ) \(\mathfrak{n} \subset \mathfrak{n}\) . Let \(z \in \mathfrak{g}\) be such that \(h = -[x, z]\) . For any integer \(n \geq 0\) , put \(\mathrm{M}_n = (\mathrm{ad}~x)^n\mathfrak{g}\) . If \(n > 0\) , we have (§1, no. 1, Lemma 1)

\[
[ \mathrm{ad} z, (\mathrm{ad} x) ^ {n} ] = n ((\mathrm{ad} h) - n + 1) (\mathrm{ad} x) ^ {n - 1}
\]

so, if \(u \in M_{n-1}\) ,

\[
n ((\operatorname{ad} h) - n + 1) u \in (\operatorname{ad} z) (\operatorname{ad} x) u + \mathrm{M} _ {n}.
\]

Since (ad \(h)\mathfrak{n}\subset \mathfrak{n}\) , it follows that

\[
((\operatorname{ad} h) - n + 1) (\mathfrak {n} \cap \mathrm{M} _ {n - 1}) \subset \mathfrak {n} \cap \mathrm{M} _ {n}.
\]

Since ad x is nilpotent, \(M_{n} = 0\) for sufficiently large n.~Consequently, the eigenvalues of ad \(h|n\) are integers \(\geq 0\) . Thus, the restriction of ad \(h + 2\) to n is invertible.

Now \([h,z] + 2z\in \mathfrak{n}\) since

\[
\begin{array}{r} [ x, [ h, z ] + 2 z ] = [ [ x, h ], z ] + [ h, [ x, z ] ] + 2 [ x, z ] \\ = [ - 2 x, z ] + [ h, - h ] + 2 [ x, z ] = 0. \end{array}
\]

Hence there exists \(z' \in \mathfrak{n}\) such that \([h, z'] + 2z' = [h, z] + 2z\) , that is, \([h, y] = -2y\) , putting \(y = z - z'\) . Since \([x, y] = [x, z] = -h\) , this completes the proof.

PROPOSITION 2 (Jacobson-Morozov). Assume that \(\mathfrak{g}\) is semi-simple. Let \(x\) be a non-zero nilpotent element of \(\mathfrak{g}\) . There exist \(h, y \in \mathfrak{g}\) such that \((x, h, y)\) is an \(\mathfrak{sl}_2\) -triplet.

Let \(\mathfrak{n}=\operatorname{Ker}(\operatorname{ad}x)^{2}\) . If \(z\in\mathfrak{n}\) , then {[}ad x, {[}ad x, ad z{]}{]} = ad({[}x, {[}x, z{]}{]}) = 0. By Lemma 5, ad \(x\circ ad\) z is nilpotent, so \(\operatorname{Tr}(\operatorname{ad}x\circ\operatorname{ad}z)=0\) . This shows that x is orthogonal to n with respect to the Killing form \(\Phi\) of g. Since

\[
\Phi ((\operatorname{ad} x) ^ {2} y, y ^ {\prime}) = \Phi (y, (\operatorname{ad} x) ^ {2} y ^ {\prime})
\]

for all \(y, y' \in g\) , and since \(\Phi\) is non-degenerate, the orthogonal complement of n is the image of \((\operatorname{ad} x)^{2}\) . Hence there exists \(y' \in g\) such that \(x = (\operatorname{ad} x)^{2}y'\) . Put

\[
h = - 2 [ x, y ^ {\prime} ];
\]

we have \([h,x]=2x\) and \(h\in(\operatorname{ad}x)(\mathfrak{g})\) . It now suffices to apply Lemma 6.

COROLLARY. Assume that \(\mathfrak{g}\) is semi-simple. Let \(\mathrm{G}\) be a group of automorphisms of \(\mathfrak{g}\) containing \(\operatorname{Aut}_e(\mathfrak{g})\) . The map which associates to any \(\mathfrak{sl}_2\) -triplet \((x,h,y)\) in g the nilpotent element x defines, by passage to the quotient, a bijection from the set of G-conjugacy classes of \(sl_{2}\) -triplets to the set of G-conjugacy classes of non-zero nilpotent elements.

This follows from Prop. 1 and 2.

Lemma 7. Let K be a commutative field with at least 4 elements. Let G be the group of matrices \(\begin{pmatrix}\alpha&\beta\\0&\alpha^{-1}\end{pmatrix}\) where \(\alpha\in K^{*},\beta\in K\) . Let \(G'\) be the group of such matrices such that \(\alpha=1\) . Then \(G'=(G,G)\) .

If \(\alpha, \alpha' \in \mathrm{K}^*\) and \(\beta, \beta' \in \mathrm{K}\) ,

\[
\begin{array}{c} \left( \begin{array}{c c} \alpha & \beta \\ 0 & \alpha^ {- 1} \end{array} \right) \left( \begin{array}{c c} \alpha^ {\prime} & \beta^ {\prime} \\ 0 & \alpha^ {- 1} \end{array} \right) \left( \begin{array}{c c} \alpha & \beta \\ 0 & \alpha^ {- 1} \end{array} \right) ^ {- 1} \left( \begin{array}{c c} \alpha^ {\prime} & \beta^ {\prime} \\ 0 & \alpha^ {- 1} \end{array} \right) ^ {- 1} \\ = \left( \begin{array}{c c} 1 & - \alpha^ {\prime} \beta^ {\prime} - \alpha \beta \alpha^ {2} + \alpha^ {2} \alpha^ {\prime} \beta^ {\prime} + \alpha \beta \\ 0 & 1 \end{array} \right). \end{array}
\]

In particular,

\[
\left( \begin{array}{c c} 1 & \beta \\ 0 & 1 \end{array} \right) \left( \begin{array}{c c} \alpha^ {\prime} & 0 \\ 0 & \alpha^ {- 1} \end{array} \right) \left( \begin{array}{c c} 1 & \beta \\ 0 & 1 \end{array} \right) ^ {- 1} \left( \begin{array}{c c} \alpha^ {\prime} & 0 \\ 0 & \alpha^ {- 1} \end{array} \right) ^ {- 1} = \left( \begin{array}{c c} 1 & \beta (1 - \alpha^ {2}) \\ 0 & 1 \end{array} \right).
\]

But there exists \(\alpha_0' \in \mathrm{K}^*\) such that \(\alpha_0' \neq 1\) and \(\alpha_0' \neq -1\) , and then \(k.(1 - \alpha_0'^2) = k\) , hence the lemma.

PROPOSITION 3. Assume that \(\mathfrak{g}\) is semi-simple. The group \(\mathrm{Aut}_e(\mathfrak{g})\) is equal to its derived group. If \(\mathfrak{g}\) is splittable, \(\mathrm{Aut}_e(\mathfrak{g})\) is the derived group of \(\mathrm{Aut}_0(\mathfrak{g})\) .

Let \(x\) be a non-zero nilpotent element of \(\mathfrak{g}\) . Choose \(h, y \in \mathfrak{g}\) be such that \((x, h, y)\) is an \(\mathfrak{sl}_2\) -triplet (Prop. 2). The subalgebra \(\mathfrak{s}\) of \(\mathfrak{g}\) generated by \((x, h, y)\) can be identified with \(\mathfrak{sl}(2, k)\) . Let \(\rho\) be the representation \(z \mapsto \operatorname{ad}_{\mathfrak{g}} z\) of \(\mathfrak{s} = \mathfrak{sl}(2, k)\) on \(\mathfrak{g}\) , and let \(\pi\) be the representation of \(\mathbf{SL}(2, k)\) compatible with \(\rho\) ( \(\S 1\) , no. 4). The image of \(\pi\) is generated by the \(\exp(t\operatorname{ad}_{\mathfrak{g}} x)\) and the \(\exp(t\operatorname{ad}_{\mathfrak{g}} y)\) with \(t \in k\) (Algebra, Chap. III, \(\S 8\) , no. 9, Prop. 17), hence is contained in \(\operatorname{Aut}_e(\mathfrak{g})\) . Since \(\mathbf{SL}(2, k)\) is equal to its derived group (Lemma 7 and loc. cit.), \(\exp(\operatorname{ad}_{\mathfrak{g}} x)\) belongs to the derived group G of \(\operatorname{Aut}_e(\mathfrak{g})\) . Hence \(\operatorname{Aut}_e(\mathfrak{g})\) is equal to G. Assume now that \(\mathfrak{g}\) is splittable. Since \(\operatorname{Aut}_0(\mathfrak{g}) / \operatorname{Aut}_e(\mathfrak{g})\) is commutative ( \(\S 5\) , no. 3, Remark 3), the preceding proves that the derived group of \(\operatorname{Aut}_0(\mathfrak{g})\) is \(\operatorname{Aut}_e(\mathfrak{g})\) .

\subsection{SIMPLE ELEMENTS}

DEFINITION 2. An element \(h\) of \(\mathfrak{g}\) is said to be simple if there exist \(x, y \in \mathfrak{g}\) such that \((x, h, y)\) is an \(\mathfrak{sl}_2\) -triplet in \(\mathfrak{g}\) .

We also say that \(h\) is the simple element of the \(\mathfrak{sl}_2\) -triplet \((x, h, y)\) .

PROPOSITION 4. Let \(h\) be a non-zero element of \(\mathfrak{g}\) . Then \(h\) is simple if and only if there exists \(x \in \mathfrak{g}\) such that \([h, x] = 2x\) and \(h \in (\operatorname{ad} x)(\mathfrak{g})\) .

The condition is clearly necessary. It is sufficient by Lemma 6.

PROPOSITION 5. Assume that g is splittable semi-simple. Let h be a splitting Cartan subalgebra of g, R the set of roots of \((\mathfrak{g},\mathfrak{h})\) , and B a basis of R. Let h be a simple element of g belonging to h. Then h is conjugate under \(\operatorname{Aut}_{e}(\mathfrak{g},\mathfrak{h})\) to an element \(h'\) of h such that \(\alpha(h') \in \{0,1,2\}\) for all \(\alpha \in B\) .

The eigenvalues of \(\mathrm{ad}_{\mathfrak{g}}h\) belong to \(\mathbf{Z}\) (§1, no. 2, Cor. of Prop. 2). Hence \(h \in \mathfrak{h}_{\mathbf{Q}}\) . There exists an element \(w\) of the Weyl group of \((\mathfrak{g}, \mathfrak{h})\) such that \(\alpha(wh) \geq 0\) for all \(\alpha \in \mathrm{B}\) (Chap. VI, §1, no. 5, Th. 2 (i)). In view of §2, no. 2, Cor. of Th. 2, we are reduced to the case in which \(\alpha(h) \in \mathbf{N}\) for all \(\alpha \in \mathrm{B}\) . Let \(\mathrm{R}_{+}\) be the set of positive roots relative to B, and \(\mathrm{R}_{-} = -\mathrm{R}_{+}\) . There exists an \(\mathfrak{sl}_2\) -triplet in \(\mathfrak{g}\) of the form \((x, h, y)\) . Let T be the set of roots \(\beta\) such that \(\beta(h) = -2\) . Then \(\mathrm{T} \subset \mathrm{R}_{-}\) and \(y \in \sum_{\beta \in \mathrm{T}} \mathfrak{g}^{\beta}\) . Assume that there exists \(\alpha \in \mathrm{B}\) such that \(\alpha(h) > 2\) . For all \(\beta \in \mathrm{T}\) , we have \((\alpha + \beta)(h) > 0\) , so \(\alpha + \beta \notin \mathrm{R}_{-}\) and \(\alpha + \beta \neq 0\) ; on the other hand, since \(\beta \in \mathrm{R}_{-}\) and \(\alpha \in \mathrm{B}\) , we have \(\alpha + \beta \notin \mathrm{R}_{+}\) ; hence \(\alpha + \beta \notin \mathrm{R} \cup \{0\}\) , so \([\mathfrak{g}^{\alpha}, \mathfrak{g}^{\beta}] = 0\) . Thus, \([y, \mathfrak{g}^{\alpha}] = 0\) . But \(\mathrm{ad}_{\mathfrak{g}}y|\mathfrak{g}^{\alpha}\) is injective since \(\alpha(h) > 0\) (§1, no. 2, Cor. of Prop. 2). This contradiction proves that \(\alpha(h) \leq 2\) for all \(\alpha \in \mathrm{B}\) .

COROLLARY. If \(k\) is algebraically closed and if \(\mathfrak{g}\) is semi-simple of rank \(l\) , the number of conjugacy classes of simple elements of \(\mathfrak{g}\) , relative to \(\operatorname{Aut}_e(\mathfrak{g})\) , is at most \(3^l\) .

Indeed, every semi-simple element of \(\mathfrak{g}\) is conjugate under \(\mathrm{Aut}_e(\mathfrak{g})\) to an element of \(\mathfrak{h}\) .

Lemma 8. Assume that \(k\) is algebraically closed and that \(\mathfrak{g}\) is semi-simple. Let \(h\) be a semi-simple element of \(\mathfrak{g}\) such that the eigenvalues of \(\operatorname{ad} h\) are rational. Let \(\mathfrak{g}^0 = \operatorname{Ker}(\operatorname{ad} h)\) , \(\mathfrak{g}^2 = \operatorname{Ker}(\operatorname{ad} h - 2)\) . Let \(G_h\) be the set of elementary automorphisms of \(\mathfrak{g}\) leaving \(h\) fixed. Let \(x \in \mathfrak{g}^2\) be such that \([x, \mathfrak{g}^0] = \mathfrak{g}^2\) . Then \(G_h x\) contains a subset of \(\mathfrak{g}^2\) that is dense and open in the Zariski topology.

Let \(\mathfrak{h}\) be a Cartan subalgebra of \(\mathfrak{g}^0\) . This is a Cartan subalgebra of \(\mathfrak{g}\) containing \(h\) (Chap. VII, §2, no. 3, Prop. 10). We have \(h \in \mathfrak{h}_{\mathbf{Q}}\) . Let R be the root system of \((\mathfrak{g}, \mathfrak{h})\) , Q the group of radical weights. There exists a basis B of R such that \(\alpha(h) \geq 0\) for all \(\alpha \in B\) .

Let \(\mathrm{U}\) be the set of \(z \in \mathfrak{h}\) such that \(\alpha(z) \neq 0\) for all \(\alpha \in \mathrm{B}\) . Let \((H_{\alpha}')_{\alpha \in \mathrm{B}}\) be the basis of \(\mathfrak{h}\) dual to \(\mathrm{B}\) . If \(z \in \mathrm{U}\) , there exists a homomorphism from \(\mathrm{Q}\) to \(k^*\) that takes any \(\gamma \in \mathrm{Q}\) to \(\prod_{\alpha \in \mathrm{B}} \alpha(z)^{\gamma(H_{\alpha}')}\) . By §5, Prop. 2 and 4, the endomorphism \(\varphi(z)\) of the vector space \(\mathfrak{g}\) which induces on \(\mathfrak{g}^{\gamma}\) the homothety with ratio \(\prod_{\alpha \in \mathrm{B}} \alpha(z)^{\gamma(H_{\alpha}')}\) is an elementary automorphism of \(\mathfrak{g}\) , which clearly belongs to \(G_h\) .

Let \(s \in \mathfrak{h}\) . If \(\gamma \in \mathrm{R}\) is such that \(\mathfrak{g}^{\gamma} \cap \mathfrak{g}^{2} \neq 0\) ,

\[
2 = \gamma (h) = \gamma \left(\sum_ {\alpha \in \mathrm{B}} \alpha (h) H _ {\alpha} ^ {\prime}\right) = \sum_ {\alpha \in \mathrm{B}} \alpha (h) \gamma (H _ {\alpha} ^ {\prime});
\]

since \(\alpha(h) \geq 0\) for all \(\alpha \in \mathrm{B}\) , and since the \(\gamma(H_{\alpha}^{\prime})\) are integers either all \(\geq 0\) or all \(\leq 0\) , we have \(\gamma(H_{\alpha}^{\prime}) \in \mathbf{N}\) for all \(\alpha \in \mathrm{B}\) . Thus, we can consider (for \(z \in \mathfrak{h}\) ) the endomorphism \(\psi(z)\) of the vector space \(\mathfrak{g}^2\) that induces on \(\mathfrak{g}^{\gamma} \cap \mathfrak{g}^2\) the homothety with ratio \(\prod_{\alpha \in \mathrm{B}} \alpha(z)^{\gamma(H_{\alpha}^{\prime})}\) . The map \(z \mapsto \psi(z)\) from \(\mathfrak{h}\) to \(\operatorname{End}(\mathfrak{g}^2)\) is polynomial. For \(z \in \mathrm{U}\) , we have \(\psi(z) = \varphi(z)|\mathfrak{g}^2\) .

Let \(\gamma_1, \ldots, \gamma_r\) be the distinct roots of \((\mathfrak{g}, \mathfrak{h})\) vanishing on \(h\) . If \(y_1 \in \mathfrak{g}^{\gamma_1}, \ldots, y_r \in \mathfrak{g}^{\gamma_r}\) , we have \(e^{\mathrm{ad} y_1} \ldots e^{\mathrm{ad} y_r} \in \mathrm{G}_h\) . We can thus define a map \(\rho\) from \(\mathfrak{h} \times \mathfrak{g}^{\gamma_1} \times \cdots \times \mathfrak{g}^{\gamma_r}\) to \(\mathfrak{g}^2\) by putting

\[
\rho (z, y _ {1}, \dots , y _ {r}) = \psi (z) e ^ {\operatorname{ad} y _ {1}} \dots e ^ {\operatorname{ad} y _ {r}} x
\]

for \(z \in \mathfrak{h}\) , \(y_1 \in \mathfrak{g}^{\gamma_1}, \ldots, y_r \in \mathfrak{g}^{\gamma_r}\) . This map is polynomial, and \(\rho(\mathrm{U}, \mathfrak{g}^{\gamma_1}, \ldots, \mathfrak{g}^{\gamma_r}) \subset \mathrm{G}_h x\) . By Chap. VII, App. I, Prop. 3 and 4, it suffices to prove that the tangent linear map of \(\rho\) is surjective at some point.

Now let T be the tangent linear map of \(z \mapsto \psi(z)\) at \(h_{0} = \sum_{\alpha \in B} H_{\alpha}'\) . Then \(\mathrm{T}(z)\) is the endomorphism of \(g^{2}\) that induces on \(g^{\gamma} \cap g^{2}\) the homothety with ratio

\[
\sum_ {\alpha \in \mathrm{B}} \gamma (H _ {\alpha} ^ {\prime}) \alpha (h _ {0}) ^ {\gamma (H _ {\alpha} ^ {\prime}) - 1} \alpha (z) \prod_ {\beta \in \mathrm{B}, \beta \neq \alpha} \beta (h _ {0}) ^ {\gamma (H _ {\alpha} ^ {\prime})} = \sum_ {\alpha \in \mathrm{B}} \gamma (H _ {\alpha} ^ {\prime}) \alpha (z) = \gamma (z).
\]

Thus, the tangent linear map of \(z \mapsto \rho(z,0,\ldots,0)\) at \(h_{0}\) is the map \(z \mapsto [z,x]\) ; its image is \([x,h]\) . The tangent linear map at 0 of the map \(y_{1} \mapsto \rho(h_{0},y_{1},0,\ldots,0)\) is the map \(y_{1} \mapsto \psi(h_{0})[y_{1},x]\) ; this last map has image \(\psi(h_{0})[x,\mathfrak{g}^{\gamma_{1}}] = [x,\mathfrak{g}^{\gamma_{1}}]\) . Similarly, the tangent linear map at 0 of the map \(y_{i} \mapsto \rho(h_{0},0,\ldots,0,y_{i},0,\ldots,0)\) has image \([x,\mathfrak{g}^{\gamma_{i}}]\) . Finally, the tangent linear map of \(\rho\) at \((h_{0},0,\ldots,0)\) has image

\[
[ x, \mathfrak {h} + \mathfrak {g} ^ {\gamma_ {1}} + \dots + \mathfrak {g} ^ {\gamma_ {r}} ] = [ x, \mathfrak {g} ^ {0} ] = \mathfrak {g} ^ {2}.\tag{Q.E.D.}
\]

*The group \(\mathrm{G}_h\) is an algebraic group with Lie algebra ad \(\mathfrak{g}^0_*\)

PROPOSITION 6. Assume that \(k\) is algebraically closed and that \(\mathfrak{g}\) is semi-simple. Let \(G\) be a group of automorphisms of \(\mathfrak{g}\) containing \(\operatorname{Aut}_e(\mathfrak{g})\) . Let \((x, h, y)\) and \((x', h', y')\) be \(\mathfrak{sl}_2\) -triplets in \(\mathfrak{g}\) . The following conditions are equivalent:

\begin{enumerate}
\def\labelenumi{(\roman{enumi})}
\item
  h and \(h'\) are G-conjugate;
\item
  \((x,h,y)\) and \((x^{\prime},h^{\prime},y^{\prime})\) are G-conjugate.
\end{enumerate}

We only have to prove the implication (i) \(\Longrightarrow\) (ii), and we are reduced immediately to the case in which \(h = h'\) . Introduce \(g^{2}\) and \(G_{h}\) as in Lemma 8. We have \(x \in g^{2}\) , and \([x, g^{0}] = g^{2}\) by §1, no. 2, Cor. of Prop. 2. Hence \(G_{h}x\) contains a subset of \(g^{2}\) that is dense and open in the Zariski topology, and so does \(G_{h}x'\) . So there exists \(a \in G_{h}\) such that \(a(x) = x'\) . We have \(a(h) = h\) , and consequently \(a(y) = y'\) (no. 1, Lemma 1).

COROLLARY 1. The map which associates to any \(\mathfrak{sl}_2\) -triplet its simple element defines by passage to the quotients a bijection from the set of G-conjugacy classes of \(\mathfrak{sl}_2\) -triplets to the set of G-conjugacy classes of simple elements.

COROLLARY 2. If \(\operatorname{rk}(\mathfrak{g}) = l\) , the number of conjugacy classes, relative to \(\operatorname{Aut}_e(\mathfrak{g})\) , of non-zero nilpotent elements of \(\mathfrak{g}\) is at most \(3^l\) .

This follows from Cor. 1, the Cor. of Prop. 2, and the Cor. of Prop. 5.

COROLLARY 3. If \(\mathrm{rk}(\mathfrak{g}) = l\) , the number of conjugacy classes, relative to \(\mathrm{Aut}_e(\mathfrak{g})\) , of subalgebras of \(\mathfrak{g}\) isomorphic to \(\mathfrak{sl}(2,k)\) is at most \(3^l\) .

This follows from Cor. 1, Prop. 1, and the Cor. of Prop. 5.

\subsection{PRINCIPAL ELEMENTS}

DEFINITION 3. Assume that g is semi-simple.

\begin{enumerate}
\def\labelenumi{(\roman{enumi})}
\item
  A nilpotent element \(x\) of \(\mathfrak{g}\) is said to be principal if the dimension of \(\operatorname{Ker} \operatorname{ad} x\) is the rank of \(\mathfrak{g}\) .
\item
  A simple element \(h\) of \(\mathfrak{g}\) is said to be principal if \(h\) is regular and the eigenvalues of \(\operatorname{ad} h\) in an algebraic closure of \(k\) belong to 2Z.
\item
  An \(\mathfrak{sl}_2\) -triplet \((x,h,y)\) of \(\mathfrak{g}\) is said to be principal if the length of \(\mathfrak{g}\) , considered as a module over \(kx + kh + ky\) , is equal to the rank of \(\mathfrak{g}\) .
\end{enumerate}

PROPOSITION 7. Assume that g is semi-simple. Let \((x,h,y)\) be an \(sl_{2}\) -triplet in g. The following conditions are equivalent:

\begin{enumerate}
\def\labelenumi{(\roman{enumi})}
\item
  \(x\) is principal;
\item
  h is principal;
\item
  \((x,h,y)\) is principal.
\end{enumerate}

For \(p \in Z\) , let \(\mathfrak{g}^{p} = \operatorname{Ker}(\operatorname{ad} h - p)\) . Let \(g' = \sum_{p \in Z} g^{2p}\) . If g is considered as a module over \(a = kx + kh + ky\) , \(g'\) is the sum of the simple submodules of odd dimension (§1, no. 2, Cor. of Prop. 2). Let l (resp. \(l'\) ) be the length of g (resp. \(g'\) ) considered as an a-module. By §1, no. 2,

\[
\dim (\operatorname{Ker} \operatorname{ad} x) = l \geq l ^ {\prime} = \dim (\operatorname{Ker} \operatorname{ad} h) \geq \operatorname{rk} (\mathfrak {g}).
\]

The equivalence of (i) and (iii) follows immediately. On the other hand, condition (ii) means that \(\dim(\operatorname{Ker} \operatorname{ad} h) = \operatorname{rk}(\mathfrak{g})\) and \(g' = g\) , in other words that

\[
\dim (\operatorname{Ker} \operatorname{ad} h) = \operatorname{rk} (\mathfrak {g})
\]

and \(l = l'\) . The equivalence of (ii) and the other conditions follows.

PROPOSITION 8. Assume that g is semi-simple \(\neq 0\) . Let h be a splitting Cartan subalgebra of g, R the root system of \((\mathfrak{g},\mathfrak{h})\) , B a basis of R, \(h^{0}\) the element of h such that \(\alpha(h^{0}) = 2\) for all \(\alpha \in B\) .

\begin{enumerate}
\def\labelenumi{(\roman{enumi})}
\item
  The element \(h^0\) is simple and principal.
\item
  The elements \(x\) of \(\mathfrak{g}\) such that there exists an \(\mathfrak{sl}_2\) -triplet of the form \((x, h^0, y)\) are the elements of \(\sum_{\alpha \in B} \mathfrak{g}^\alpha\) that have a non-zero component in each \(\mathfrak{g}^\alpha\) .
\end{enumerate}

The element \(h^0\) is that considered in §7, no. 5, Lemma 2 (cf.~loc. cit., formula (1)). It follows from this lemma that \(h^0\) is simple principal and that, if \(x \in \sum_{\alpha \in \mathrm{B}} \mathfrak{g}^\alpha\) has a non-zero component in each \(\mathfrak{g}^\alpha\) , there exists an \(\mathfrak{sl}_2\) -triplet of the form \((x, h^0, y)\) . Conversely, let \((x, h^0, y)\) be a \(\mathfrak{sl}_2\) -triplet. We have \([h^0, x] = 2x\) , so \(x \in \sum_{\gamma \in \mathrm{R}, \gamma(h^0) = 2} \mathfrak{g}^\gamma = \sum_{\alpha \in \mathrm{B}} \mathfrak{g}^\alpha\) . Similarly, \(y \in \sum_{\alpha \in \mathrm{B}} \mathfrak{g}^{-\alpha}\) . Write

\[
h ^ {0} = \sum_ {\alpha \in \mathrm{B}} a _ {\alpha} H _ {\alpha} \quad \text { where } a _ {\alpha} > 0 \text { for all } \alpha \in \mathrm{B},
\]

\[
x = \sum_ {\alpha \in \mathrm{B}} X _ {\alpha} \qquad \text { where } X _ {\alpha} \in \mathfrak {g} ^ {\alpha} \text { for   all } \alpha \in \mathrm{B},
\]

\[
y = \sum_ {\alpha \in \mathrm{B}} X _ {- \alpha} \quad \text { where } X _ {- \alpha} \in \mathfrak {g} ^ {- \alpha} \text { for   all } \alpha \in \mathrm{B}.
\]

Then

\[
\sum_ {\alpha \in \mathrm{B}} a _ {\alpha} H _ {\alpha} = h ^ {0} = [ y, x ] = \sum_ {\alpha , \beta \in \mathrm{B}} [ X _ {- \beta}, X _ {\alpha} ] = \sum_ {\alpha \in \mathrm{B}} [ X _ {- \alpha}, X _ {\alpha} ]
\]

so \([X_{-\alpha}, X_{\alpha}] \neq 0\) for all \(\alpha \in \mathrm{B}\) .

COROLLARY. In a splittable semi-simple Lie algebra, there exist principal nilpotent elements.

In a non-splittable semi-simple Lie algebra, 0 may be the only nilpotent element.

PROPOSITION 9. Assume that \(k\) is algebraically closed and that \(\mathfrak{g}\) is semi-simple. All the principal simple (resp. nilpotent) elements of \(\mathfrak{g}\) are conjugate under \(\operatorname{Aut}_e(\mathfrak{g})\) .

We retain the notations of Prop. 8. Let h be a principal simple element. It is conjugate under \(\operatorname{Aut}_{e}(\mathfrak{g})\) to an \(h' \in h\) such that \(\alpha(h') \in \{0,1,2\}\) for all \(\alpha \in B\) (no. 3, Prop. 5). Since \(h'\) is principal simple, \(\alpha(h') \neq 0\) and \(\alpha(h') \in 2Z\) for all \(\alpha \in B\) , so \(\alpha(h') = 2\) for all \(\alpha \in B\) , and hence \(h' = h^{0}\) . This proves the assertion for principal simple elements.

Let \(x, x'\) be principal nilpotent elements. There exist \(\mathfrak{sl}_2\) -triplets \((x, h, y)\) , \((x', h', y')\) . By Prop. 7, \(h\) and \(h'\) are principal simple, hence conjugate under \(\mathrm{Aut}_e(\mathfrak{g})\) by the preceding. So \(x\) and \(x'\) are conjugate under \(\mathrm{Aut}_e(\mathfrak{g})\) (Prop. 6).

Lemma 9. With the notations of Prop. 8, put \(\mathfrak{g}^p = \mathrm{Ker}(\mathrm{ad}h^0 - p)\) for \(p \in \mathbf{Z}\) . Let \(\mathfrak{g}_*^2\) be the set of elements of \(\mathfrak{g}^2 = \sum_{\alpha \in \mathrm{B}} \mathfrak{g}^\alpha\) that have a non-zero component in each \(\mathfrak{g}^\alpha\) . Let \(\mathrm{R}_+\) be the set of positive roots relative to \(\mathrm{B}\) , \(\mathfrak{n}_+ = \sum_{\alpha \in \mathrm{R}_+} \mathfrak{g}^\alpha\) , and \(x \in \mathfrak{g}_*^2\) . Then \(e^{\mathrm{ad}\mathfrak{n}_+}. x = x + [\mathfrak{n}_+, \mathfrak{n}_+]\) .

It is clear that \(e^{\mathrm{ad} \mathfrak{n}_+}.x \subset x + [\mathfrak{n}_+, \mathfrak{n}_+]\) . We prove that, if \(v \in [\mathfrak{n}_+, \mathfrak{n}_+]\) , then \(x + v \in e^{\mathrm{ad} \mathfrak{n}_+}.x\) . Put \(\mathfrak{n}^{(p)} = \sum_{r \geq p} \mathfrak{g}^{2r}\) ; it suffices to prove that

\[
x + v \in e ^ {\mathrm{ad} \mathfrak {n} _ {+}}. x + \mathfrak {n} ^ {(p)}
\]

for all \(p \geq 2\) . This is clear for \(p = 2\) since \(\mathfrak{n}^{(2)} = [\mathfrak{n}_+, \mathfrak{n}_+]\) (§3, no. 3, Prop. 9 (iii)). Assume that we have found \(z \in \mathfrak{n}_+\) such that \(v + x - e^{\mathrm{ad}z}.x \in \mathfrak{n}^{(p)}\) . Since there exists an \(\mathfrak{sl}_2\) -triplet of the form \((x, h^0, y)\) (Prop. 8), §1, no. 2, Cor. of Prop. 2 proves that \([x, \mathfrak{g}^{2p-2}] = \mathfrak{g}^{2p}\) ; hence, there exists \(z' \in \mathfrak{g}^{2p-2} \subset \mathfrak{n}_+\) such that

\[
v + x - e ^ {\mathrm{ad} z}. x \in [ z ^ {\prime}, x ] + \mathfrak {n} ^ {(p + 1)}.
\]

So \(v + x \in e^{\mathrm{ad}(z + z')}.x + \mathfrak{n}^{(p + 1)}\) , and our assertion is established by induction.

PROPOSITION 10. Assume that g is semi-simple. Let h be a splitting Cartan subalgebra of g, R the root system of \((\mathfrak{g},\mathfrak{h})\) , B a basis of R, \(R_{+}\) the set of positive roots relative to B, and \(n_{+}=\sum_{\alpha\in R_{+}}g^{\alpha}\) . The principal nilpotent elements belonging to \(n_{+}\) are the elements of \(n_{+}\) having a non-zero component in \(g^{\alpha}\) for all \(\alpha\in B\) .

Prop. 8 and Lemma 9 prove that such elements are principal nilpotent. We prove the converse. Evidently we can assume that \(\mathfrak{g}\) is simple. Let \(h^0\) and \(\mathfrak{g}^p\) be as in Prop. 8 and Lemma 9. Let \(\omega\) be the highest root, and put \(\omega(h^0) = 2q\) ; we have \(q = h - 1\) , where \(h\) is the Coxeter number of R, cf.~Chap. VI, §1, no. 11, Prop. 31. Then \(\mathfrak{g}^{2q} = \mathfrak{g}^\omega\) , \(\mathfrak{g}^{-2q} = \mathfrak{g}^{-\omega}\) , and \(\mathfrak{g}^{2k} = 0\) for \(|k| > q\) . There exists a principal \(\mathfrak{sl}_2\) -triplet \((x^0, h^0, y^0)\) . By §1, no. 2, Cor. of Prop. 2, \((\mathrm{ad} x^0)^{2q}(\mathfrak{g}^{-\omega}) = \mathfrak{g}^\omega\) , so \((\mathrm{ad} x^0)^{2q} \neq 0\) . Let \(x\) be a principal nilpotent element of \(\mathfrak{g}\) belonging to \(\mathfrak{n}_+\) . If \(\bar{k}\) is an algebraic closure of \(k\) , \(x \otimes 1\) and \(x^0 \otimes 1\) are conjugate under an automorphism of \(\mathfrak{g} \otimes_k \bar{k}\) (Prop. 9), so \((\mathrm{ad} x)^{2q} \neq 0\) . There exists \(\lambda \in \mathbb{R}\) such that \((\mathrm{ad} x)^{2q} \mathfrak{g}^\lambda \neq 0\) . Put \(x = \sum_{n \geq 1} x_n\) , where \(x_n \in \mathfrak{g}^{2n}\) . Then

\[
(\operatorname{ad} x) ^ {2 q} \mathfrak {g} ^ {\lambda} \subset (\operatorname{ad} x _ {1}) ^ {2 q} \mathfrak {g} ^ {\lambda} + \sum_ {k > 4 q + \lambda (h ^ {0})} \mathfrak {g} ^ {k} = (\operatorname{ad} x _ {1}) ^ {2 q} \mathfrak {g} ^ {\lambda},
\]

since \(4q + \lambda(h^0) \geq 4q - 2q = 2q\) . Now \((\operatorname{ad} x_1)^{2q} \mathfrak{g}^\lambda \subset \mathfrak{g}^{4q + \lambda(h^0)}\) , where \(\lambda = -\omega\) . Thus, \((\operatorname{ad} x_1)^{2q} \mathfrak{g}^{-\omega} = \mathfrak{g}^\omega\) . We have \(\omega = \sum_{\alpha \in \mathrm{B}} n_\alpha \alpha\) with \(n_\alpha > 0\) for all \(\alpha \in \mathrm{B}\) (Chap. VI, §1, no. 8, Remark). If there exists \(\alpha_0 \in \mathrm{B}\) such that \(x_1 \in \sum_{\alpha \in \mathrm{B}, \alpha \neq \alpha_0} \mathfrak{g}^\alpha\) , the relation

\[
\omega \notin - \omega + \sum_ {\alpha \in \mathrm{B}, \alpha \neq \alpha_ {0}} k \alpha
\]

implies that \(\mathfrak{g}^{\omega} \not\subset (\operatorname{ad} x_1)^p \mathfrak{g}^{-\omega}\) for all \(p\) ; this is absurd, so the component of \(x_1\) in \(\mathfrak{g}^{\alpha}\) is non-zero for all \(\alpha \in B\) .

\section{CHEVALLEY ORDERS}

\subsection{LATTICES AND ORDERS}

Let V be a Q-vector space. A lattice in V is a free Z-submodule V of V such that the Q-linear map \(\alpha_{\mathcal{V},V} : \mathcal{V} \otimes_{Z} Q \to V\) induced by the injection of V into V is bijective. When V is finite dimensional, this is the same as saying that V is a Z-submodule of finite type which generates the Q-vector space V (recall that a torsion-free Z-module of finite type is free by Algebra, Chap. VII, §4, no. 4, Cor. 2); moreover, in this case our definition is a special case of that of Commutative Algebra, Chap. VII, §4, no. 1, Def. 1 (loc. cit., Example 3). If W is a vector subspace of V, and V is a lattice in V, then \(V \cap W\) is a lattice in W.

If V is a Q-algebra, an order in V is a lattice V in the underlying vector space that is a Z-subalgebra of V; the map \(\alpha_{V,V}\) is then an isomorphism of Q-algebras. If V is a unital Q-algebra, a unital order in V is an order in V containing the unit element.

Assume that V is a Q-bigebra, with coproduct c and counit \(\gamma\) . If V is a lattice in the vector space V, the canonical map \(i: V \otimes_{Z} V \to V \otimes_{Q} V\) is injective; a biorder in V is a unital order V in the unital algebra V such that \(\gamma(\mathcal{V}) \subset \mathbf{Z}\) and \(c(\mathcal{V}) \subset i(\mathcal{V} \otimes_{\mathbf{Z}} \mathcal{V})\) ; the maps

\[
\gamma_ {\mathcal {V}} \colon \mathcal {V} \to \mathbf {Z} \text { and } c _ {\mathcal {V}} \colon \mathcal {V} \to \mathcal {V} \otimes_ {\mathbf {Z}} \mathcal {V}
\]

induced by \(\gamma\) and c give V the structure of a Z-bigebra, and the map \(\alpha_{\mathcal{V},V}\) is then an isomorphism of Q-bigebras.

\subsection{DIVIDED POWERS IN A BIGEBRA}

Let A be a unital k-algebra, \(x \in A\) , \(d \in k\) , \(n \in N\) . Put

\[
x ^ {(n, d)} = \frac {x (x - d) \dots (x - d (n - 1))}{n !} = \prod_ {i = 0} ^ {n - 1} (x - i d) / (i + 1).\tag{1}
\]

In particular, \(x^{(0,d)} = 1\) , \(x^{(1,d)} = x\) . We agree that \(x^{(n,d)} = 0\) for \(n\) an integer \(< 0\) . Put

\[
x ^ {(n)} = x ^ {(n, 0)} = \frac {x ^ {n}}{n !}\tag{2}
\]

\[
\binom {x} {n} = x ^ {(n, 1)} = \frac {x (x - 1) \ldots (x - n + 1)}{n !}.\tag{3}
\]

PROPOSITION 1. Let A be a bigebra, with coproduct \(c\) , and \(x\) a primitive element (Chap. II, §1, no. 2) of A. Then

\[
c (x ^ {(n, d)}) = \sum_ {p \in \mathbf {N}, q \in \mathbf {N}, p + q = n} x ^ {(p, d)} \otimes x ^ {(q, d)}.\tag{4}
\]

The proposition is trivial for \(n \leq 0\) . We argue by induction on \(n\) . If formula (4) is true for \(n\) , then

\[
\begin{array}{l} (n + 1) c (x ^ {(n + 1, d)}) = c (x - d n) c (x ^ {(n, d)}) \\ \quad = (x \otimes 1 + 1 \otimes x - d n   1 \otimes 1) c (x ^ {(n, d)}) \\ \quad = \sum_ {p + q = n} [ x x ^ {(p, d)} \otimes x ^ {(q, d)} + x ^ {(p, d)} \otimes x x ^ {(q, d)} - (p + q) d x ^ {(p, d)} \otimes x ^ {(q, d)} ] \\ \quad = \sum_ {p + q = n} (x - p d) x ^ {(p, d)} \otimes x ^ {(q, d)} + \sum_ {p + q = n} x ^ {(p, d)} \otimes (x - q d) x ^ {(q, d)} \\ \quad = \sum_ {p + q = n} (p + 1) x ^ {(p + 1, d)} \otimes x ^ {(q, d)} + \sum_ {p + q = n} (q + 1) x ^ {(p, d)} \otimes x ^ {(q + 1, d)} \\ \quad = \sum_ {r + s = n + 1} r x ^ {(r, d)} \otimes x ^ {(s, d)} + \sum_ {r + s = n + 1} s x ^ {(r, d)} \otimes x ^ {(s, d)} \\ \quad = (n + 1) \sum_ {r + s = n + 1} x ^ {(r, d)} \otimes x ^ {(s, d)}, \end{array}
\]

hence formula (4) for \(n + 1\) .

\subsection{INTEGRAL VARIANT OF THE POINCAR\'E-BIRKHOFF-WITT THEOREM}

Let g be a finite dimensional Q-Lie algebra, \(U(\mathfrak{g})\) its enveloping bigebra. If I is a totally ordered set, \(\mathbf{x} = (x_i)_{i \in \mathrm{I}}\) a family of elements of g, and \(\mathbf{n} = (n_i)_{i \in \mathrm{I}} \in \mathbf{N}^{(\mathrm{I})}\) a multi-index, put

\[
\mathbf {x} ^ {(\mathbf {n})} = \prod_ {i \in \mathrm{I}} \frac {x _ {i} ^ {n _ {i}}}{n _ {i} !},\tag{5}
\]

the product being calculated in \(\mathrm{U}(\mathfrak{g})\) in accordance with the ordered set I.

THEOREM 1. Let \(\mathcal{U}\) be a biorder in the bigebra \(\mathrm{U}(\mathfrak{g})\) . Let \(\mathcal{G} = \mathcal{U} \cap \mathfrak{g}\) , which is an order in the Lie algebra \(\mathfrak{g}\) . Let \((x_i)_{i \in \mathrm{I}}\) be a basis of \(\mathcal{G}\) . Give I a total order, and assume that we are given, for all \(\mathbf{n} \in \mathbf{N}^{\mathrm{I}}\) , an element {[}n{]} of \(\mathcal{U}\) such that {[}n{]} - \(\mathbf{x}^{(\mathbf{n})}\) has filtration \(< |\mathbf{n}|\) in \(\mathrm{U}(\mathfrak{g})\) . Then, the family of the {[}n{]} for \(\mathbf{n} \in \mathbf{N}^{\mathrm{I}}\) is a basis of the Z-module \(\mathcal{U}\) .

For \(p \in N\) , let \(U_{p}(\mathfrak{g})\) be the set of elements of \(U(\mathfrak{g})\) of filtration \(\leq p\) ; then the images in \(U_{p}(\mathfrak{g}) / U_{p-1}(\mathfrak{g})\) of the \(\mathbf{x}^{(\mathbf{n})}\) such that \(|n| = p\) form a basis of this Q-vector space (Chap. I, §2, no. 7, Th. 1); hence the {[}n{]} form a basis of the Q-vector space \(U(\mathfrak{g})\) . It remains to prove the following assertion (in which we put \(M = N^{I}\) ):

\((^{*})\) if \(u\in \mathcal{U}\) , \((a_{\mathbf{n}})\in \mathbf{Z}^{(\mathrm{M})}\) , and \(d\in \mathbf{N} - \{0\}\) are such that

\[
d u = \sum_ {\mathbf {n} \in \mathrm{M}} a _ {\mathbf {n}} [ \mathbf {n} ],\tag{6}
\]

then \(d\) divides each \(a_{\mathbf{n}}\) .

For each integer \(r \geq 0\) , introduce the iterated coproduct

\[
c _ {i}: \mathcal {U} \rightarrow \mathbf {T} ^ {r} (\mathcal {U}) = \mathcal {U} \otimes \mathcal {U} \otimes \dots \otimes \mathcal {U};
\]

by definition, \(c_0\) is the counit of \(\mathcal{U}\) , \(c_1 = \mathrm{Id}_{\mathcal{U}}\) , \(c_2 = c\) (the coproduct of \(\mathcal{U}\) ), and, for \(r \geq 2\) , \(c_{r+1}\) is defined as the composite \(p \circ (c_r \otimes 1) \circ c\) :

\[
\mathcal {U} \stackrel {{c}} {{\longrightarrow}} \mathcal {U} \otimes_ {\mathbf {Z}} \mathcal {U} \stackrel {{c _ {r} \otimes 1}} {{\longrightarrow}} \mathbf {T} ^ {r} (\mathcal {U}) \otimes_ {\mathbf {Z}} \mathcal {U} \stackrel {{p}} {{\longrightarrow}} \mathbf {T} ^ {r + 1} (\mathcal {U})
\]

where p is defined by using the multiplication in the algebra \(\mathbf{T}(\mathcal{U})\) . Further, consider the canonical projection \(\pi\) of U onto \(U^{+} = Ker c_{0}\) , and the composite

\[
c _ {r} ^ {+} = \mathbf {T} ^ {r} (\pi) \circ c _ {r}: \mathcal {U} \to \mathbf {T} ^ {r} (\mathcal {U} ^ {+}).
\]

Lemma 1. Let \(\mathbf{n} \in \mathbf{N}^{\mathrm{I}}\) . If \(|\mathbf{n}| < r\) , then \(c_{r}^{+}([ \mathbf{n}]) = 0\) . If \(|\mathbf{n}| = r\) , then

\[
c _ {r} ^ {+} ([ \mathbf {n} ]) = \sum_ {\varphi} x _ {\varphi (1)} \otimes x _ {\varphi (2)} \otimes \dots \otimes x _ {\varphi (r)},\tag{7}
\]

where \(\varphi\) belongs to the set of maps from \(\{1,2,\ldots,r\}\) to I which take each value \(i\in I\) \(n_{i}\) times.

By Prop. 1,

\[
c _ {r} \left(\mathbf {x} ^ {(\mathbf {n})}\right) = \sum \mathbf {x} ^ {\left(\mathbf {p} _ {1}\right)} \otimes \dots \otimes \mathbf {x} ^ {\left(\mathbf {p} _ {r}\right)}
\]

where the summation extends over the set of sequences \((\mathbf{p}_{1},\ldots,\mathbf{p}_{r})\) of r elements of M such that \(p_{1}+\cdots+p_{r}=n\) . In view of Chap. II, §1, no. 3, Prop. 6, the map \(c_{r}^{+}\) , extended by linearity to a map from \(\mathrm{U}(\mathfrak{g})\) to \(\mathbf{T}^{r}(\mathrm{U}^{+}(\mathfrak{g}))\) , vanishes on \(\mathrm{U}_{r-1}(\mathfrak{g})\) . It follows that, for \(r\geq|\mathbf{n}|\) ,

\[
c _ {r} ^ {+} ([ \mathbf {n} ]) = c _ {r} ^ {+} (\mathbf {x} ^ {(\mathbf {n})}) = \sum \pi (\mathbf {x} ^ {(\mathbf {p} _ {1})}) \otimes \dots \otimes \pi (\mathbf {x} ^ {(\mathbf {p} _ {r})}).\tag{8}
\]

For \(r > |n|\) , the relation \(p_{1} + \cdots + p_{r} = n\) implies that at least one of the \(p_{i}\) is zero, so \(c_{r}^{+}([n]) = 0\) . For \(r = |n|\) , the only non-zero terms of the third member of (8) are those for which \(|p_{1}| = \cdots = |p_{r}| = 1\) , hence (7).

We return to the proof of Th. 1. We retain the notations of (*) and prove, by descending induction on \(|n|\) , that d divides \(a_{n}\) , which is clear when \(|n|\) is sufficiently large. If d divides \(a_{n}\) for \(|n| > r\) then, putting

\[
u ^ {\prime} = u - \sum_ {| \mathbf {n} | > r} (a _ {\mathbf {n}} / d) [ \mathbf {n} ] \in \mathcal {U},
\]

we have

\[
d u ^ {\prime} = \sum_ {| \mathbf {n} | \leq r} a _ {\mathbf {n}} [ \mathbf {n} ].\tag{9}
\]

For any map \(\varphi\) from \(\{1,2,\ldots,r\}\) to I, put

\[
e _ {\varphi} = x _ {\varphi (1)} \otimes \dots \otimes x _ {\varphi (r)}
\]

and \(a_{\varphi} = a_{\mathbf{n}}\) where \(\mathbf{n} = (\operatorname{Card}\varphi^{-1}(i))_{i\in \mathrm{I}}\) . By Lemma 1, (9) implies that

\[
d c _ {r} ^ {+} (u ^ {\prime}) = \sum_ {\varphi \in \mathrm{I} ^ {r}} a _ {\varphi} e _ {\varphi}\tag{10}
\]

so \(c_r^+(u') \in \mathbf{T}^r(\mathcal{U}^+) \cap \mathbf{Q}\mathbf{T}^r(\mathcal{G})\) . But the submodule \(\mathcal{G}\) of \(\mathcal{U}^+\) is a direct factor (Algebra, Chap. VII, §4, no. 3, Cor. of Th. 1), so the submodule \(\mathbf{T}^r(\mathcal{G})\) is a direct factor of \(\mathbf{T}^r(\mathcal{U}^+)\) , and hence \(c_r^+(u') \in \mathbf{T}^r(\mathcal{G})\) . On the other hand, the \(x_i\) form a basis of \(\mathcal{G}\) by hypothesis, so the \(e_\varphi\) form a basis of \(\mathbf{T}^r(\mathcal{G})\) . Then (10) proves that \(d\) divides the \(a_\varphi\) , that is, the \(a_{\mathbf{n}}\) for \(|\mathbf{n}| = r\) . This proves (*).